
\documentclass[11pt]{article}
\usepackage[margin=1in]{geometry}

% Core packages
\usepackage{amsmath,amssymb,amsthm,mathtools}
\usepackage{tikz-cd}
\usepackage{multicol}
\usepackage[T1]{fontenc}
\usepackage{lmodern}
\usepackage{hyperref}

% Paragraphs
\setlength{\parindent}{0pt}
\setlength{\parskip}{0.7\baselineskip}

\newtheorem{definition}{Definition}
\newtheorem{theorem}{Theorem}
\newtheorem{proposition}{Proposition}
\newtheorem{corollary}{Corollary}
\newtheorem{remark}{Remark}
\newtheorem{assumption}{Assumption}

\title{Constrained Dataset Predictive Tuning:\\A Formal Algorithmic Specification}
\author{Noah Cashin}
\date{\today}

\begin{document}
\maketitle

\section{Purpose and Scope}\label{sec:purpose-scope}

This document specifies a constrained predictive/compression tuning capability for a fixed dataset.
The specification is intended to be:
\begin{itemize}
    \item mathematically precise,
    \item directly implementable in \texttt{infotheory},
    \item optimizer-agnostic at the interface layer,
    \item runtime-selectable between annealed hill climbing,
    MC-AIXI(FAC-CTW), discounted AIQI, and a distinct exact
    warm-start controller, and
    \item mathematically explicit about objective alignment in each controller path.
\end{itemize}

This capability is \emph{not} an unconstrained architecture search engine nor a replacement for normal compression execution; it is a bounded optimization wrapper over the same compression semantics.

\section{Formal Problem Statement}\label{sec:formal-problem}

\subsection{Objects}

Let:
\begin{itemize}
    \item $x \in \{0,\dots,255\}^{n}$ denote the dataset bytes loaded from \texttt{input\_path},
    \item $\mathcal{Z}$ denote the full configuration space,
    \item $B$ denote the user-provided bounds specification,
    \item $H$ denote the fixed deployability profile (hardware/evaluator methodology and resource accounting),
    \item $z_0 \in \mathcal{Z}$ denote the validated baseline configuration,
    \item $T_{\mathrm{eval}} > 0$ denote the hard per-candidate wall-clock limit,
    \item $T_{\mathrm{total}} > 0$ denotes the hard total tuning wall-clock budget,
    \item $\theta_{\min}>0$ denote the required minimum throughput (bytes/s),
    \item $\mu_{\max}>0$ denote the required memory cap (bytes).
\end{itemize}

Define the bounded-valid candidate set:
\[
\mathcal{F}(B,H):=\{z \in \mathcal{Z}:\ \mathrm{Valid}_{B,H}(z)=1\}.
\]

For each candidate $z$, evaluation returns
\[
M(z) = (\mathrm{status}(z), c(z), \ell(z), \tau(z), \mu(z), \operatorname{diag}(z)),
\]
where
\begin{itemize}
    \item $\mathrm{status}(z) \in \{\mathrm{success},\mathrm{timeout},\mathrm{invalid},\mathrm{error}\}$,
    \item $c(z) \in \mathbb{N}$ is compressed size in bytes (defined on success),
    \item $\ell(z) \in \mathbb{R}_{\ge 0}$ is cumulative negative log-loss (diagnostic),
    \item $\tau(z) \in \mathbb{R}_{\ge 0}$ is elapsed wall-clock time,
    \item $\mu(z) \in \mathbb{R}_{\ge 0}$ is measured peak memory usage,
    \item $\operatorname{diag}(z)$ is optional diagnostics (including invalid-reason tags).
\end{itemize}

\begin{remark}[Peak-memory measurement semantics]
Each candidate evaluation is executed in its own isolated evaluator instance.

Preferred definition: $\mu(z)$ is the peak memory usage attributed to that
instance by the operating-system or container memory controller (for example,
cgroup peak memory usage) over the lifetime of the evaluation.

Fallback definition when such accounting is unavailable: the evaluator must run
as a single process, and $\mu(z)$ is the operating-system reported peak RSS of
that process.

Naive summation of RSS across a multi-process tree is not normative, because it
can double-count shared pages.
\end{remark}

Timeout semantics are hard:
\[
\tau(z) \ge T_{\mathrm{eval}} \implies \mathrm{status}(z)=\mathrm{timeout}.
\]

For successful candidates, define throughput
\[
\theta(z):=
\begin{cases}
\dfrac{n}{\tau(z)}, & \mathrm{status}(z)=\mathrm{success},\ \tau(z)>0,\\[4pt]
+\infty, & \mathrm{status}(z)=\mathrm{success},\ \tau(z)=0.
\end{cases}
\]
Define the equivalent throughput runtime cap
\[
\tau_{\theta}:=\frac{n}{\theta_{\min}}.
\]
Then, for successful candidates,
\[
\theta(z)\ge\theta_{\min}
\iff
\tau(z)\le\tau_{\theta}.
\]

\subsection{Primary Objective}

Define model-code length by canonical bytes
\[
L_B(z):=\big|\operatorname{ser}(\operatorname{canon}(z))\big|.
\]

For successful candidates, define dataset code length
\[
C_x(z):=c(z).
\]

Define the deployable set
\[
\mathcal{F}^{\mathrm{dep}}(B,H):=
\left\{z\in\mathcal{F}(B,H):\ \mathrm{status}(z)=\mathrm{success}\ \land\ \theta(z)\ge\theta_{\min}\ \land\ \mu(z)\le\mu_{\max}\right\}.
\]

The normative objective is the totalized deployable two-part MDL criterion
\[
\mathcal{J}_H(z;x)=
\begin{cases}
L_B(z)+C_x(z), & z\in\mathcal{F}^{\mathrm{dep}}(B,H),\\[3pt]
+\infty, & \text{otherwise.}
\end{cases}
\]

The primary optimization problem is
\[
\min_{z\in\mathcal{F}(B,H)}\mathcal{J}_H(z;x).
\]

\begin{remark}[Operational length units and reporting]
Operational length quantities in this specification are stored and reported in
bytes: $c(z)$, $L_B(z)$, $C_x(z)$, and $\mathcal{J}_H(z;x)$.
Bit-normalized quantities are introduced only in
Section~\ref{sec:restricted-mixture} for semantic/theoretical statements via
the conversion $\text{bits}=8\times\text{bytes}$.
\end{remark}

\subsection{Deterministic Incumbent Key and Output Semantics}

For each evaluated candidate $z$, define the deterministic key
\[
K(z):=\big(\mathcal{J}_H(z;x),\operatorname{ser}(\operatorname{canon}(z))\big),
\]
ordered lexicographically, where $\operatorname{canon}:\mathcal{Z}\to\mathcal{Z}$
and $\operatorname{ser}:\mathcal{Z}\to\{0,\dots,255\}^*$ are defined in
Section~\ref{sec:canon-cache}.

For byte strings, use standard lexicographic byte order with strict-prefix
rule: $u\prec_{\mathrm{bytes}} v$ if either (i) at the first differing byte,
$u$ has the smaller byte, or (ii) $u$ is a strict prefix of $v$.

Let
\[
\mathcal{S}_t:=\left\{z:\ z\text{ has been evaluated by step }t\ \text{and }z\in\mathcal{F}^{\mathrm{dep}}(B,H)\right\}.
\]
The incumbent is
\[
b_t:=\arg\min_{z\in\mathcal{S}_t}K(z).
\]
Write $\prec_K$ for lexicographic order on key tuples, i.e.
$K(z_1)\prec_K K(z_2)$ iff the tuple of $z_1$ is lexicographically smaller than that of $z_2$.

\begin{proposition}[Uniqueness of incumbent under finite deployable evaluated set]
If $\mathcal{S}_t$ is finite and nonempty, then $b_t$ exists and is unique.
\end{proposition}

\begin{proof}
By construction, $\prec_{\mathrm{bytes}}$ is a total order on
$\{0,\dots,255\}^*$. Hence lexicographic order on
$\big(\mathbb{R}_{\ge 0}\sqcup\{+\infty\}\big)\times\{0,\dots,255\}^*$ is total. A finite
nonempty totally ordered set has a unique minimum.
\end{proof}

This key preserves the primary optimization semantics: $\mathcal{J}_H$ is
primary, and deterministic canonical bytes are a tie-break only.

\section{Input Contract and Defaults}\label{sec:input-contract}

\subsection{Required fields}

\begin{itemize}
    \item \texttt{input\_path: Path}
    \item \texttt{baseline\_config\_path: Path}
    \item \texttt{algorithm: ControllerKind}
    \item \texttt{bounds: BoundsSpec}
    \item \texttt{eval\_time\_limit\_seconds: f64}
    \item \texttt{time\_budget\_seconds: f64}
    \item \texttt{min\_throughput\_bytes\_per\_second: f64} with constraint $>0$
    \item \texttt{max\_memory\_bytes: u64} with constraint $>0$
    \item \texttt{output\_config\_path: Path}
    \item \texttt{seed: u64}
\end{itemize}

Define the throughput floor parameter by
\[
\theta_{\min}:=\texttt{min\_throughput\_bytes\_per\_second}\in\mathbb{R}_{>0}.
\]

Define the memory-cap parameter by
\[
\mu_{\max}:=\texttt{max\_memory\_bytes}\in\mathbb{R}_{>0}.
\]

Define runtime controller kind by
\[
\texttt{ControllerKind}\in
\left\{
\begin{aligned}
&\texttt{"annealed\_hill\_climbing"},\ \texttt{"mc\_aixi\_fac\_ctw"},\\
&\texttt{"aiqi\_discounted"},\ \texttt{"aiqi\_warmstart\_exact\_jh"}
\end{aligned}
\right\}.
\]

Runtime dispatch is normative:
\begin{itemize}
    \item if \texttt{algorithm = "annealed\_hill\_climbing"}, execute
    Section~\ref{sec:ahc-runtime},
    \item if \texttt{algorithm} is \texttt{mc\_aixi\_fac\_ctw},
    \texttt{aiqi\_discounted}, or
    \texttt{aiqi\_warmstart\_exact\allowbreak\_jh},
    execute
    Section~\ref{sec:aixi-runtime}.
\end{itemize}

Controller-specific required fields:
\begin{itemize}
    \item for \texttt{mc\_aixi\_fac\_ctw}, \texttt{aiqi\_discounted}, and
    \texttt{aiqi\_warmstart\_exact\allowbreak\_jh}:
    \begin{itemize}
        \item \texttt{agent\_actions}
        \item \texttt{observation\_bits}
        \item \texttt{observation\_stream\_len}
        \item \texttt{reward\_bits}
        \item \texttt{min\_reward}
        \item \texttt{max\_reward}
        \item \texttt{reward\_offset}
        \item \texttt{planner\_simulations\_per\_step}\((\in\mathbb{N},\ge 1)\)
    \end{itemize}
    \item for \texttt{aiqi\_discounted} additionally:
    \begin{itemize}
        \item \texttt{return\_horizon}\(\in\mathbb{N},\ge 1\)
        \item \texttt{return\_bins}\(\in\mathbb{N},\text{power of two}\)
        \item \texttt{discount\_factor}\(\in[0,1)\)
    \end{itemize}
    \item for \texttt{aiqi\_warmstart\_exact\allowbreak\_jh} additionally:
    \begin{itemize}
        \item \texttt{return\_horizon}\(\in\mathbb{N},\ge 1\)
        \item \texttt{warmstart\_teacher\_dataset\_path: Path}
        (initial same-task teacher data)
        \item \texttt{label\_phase\_period}
        \((\in\mathbb{N},\ge\texttt{return\_horizon})\)
    \end{itemize}
    \item for \texttt{"annealed\_hill\_climbing"}:
    \texttt{max\_mutation\_radius}.
\end{itemize}

\subsection{Normative input precondition: deployable baseline}\label{subsec:baseline-precondition}

After loading and canonicalizing the baseline candidate
\[
z_0^{\mathrm{can}}:=\operatorname{canon}(z_0),
\]
the runtime must perform the normative baseline evaluation before entering any
controller loop and require
\[
\mathrm{status}(z_0^{\mathrm{can}})=\mathrm{success}
\quad\land\quad
\theta(z_0^{\mathrm{can}})\ge\theta_{\min}
\quad\land\quad
\mu(z_0^{\mathrm{can}})\le\mu_{\max}.
\]
If this precondition fails, execution must abort with no optimization steps
performed. This precondition is normative for all controller choices.

\subsection{Execution controls}

Recommended controls include:
\begin{itemize}
    \item \texttt{max\_evaluations}
    \item \texttt{cpu\_affinity}
    \item \texttt{threads}
    \item \texttt{warmup\_baseline\_runs}
    \item \texttt{self\_improvement\_rounds}
    \item \texttt{max\_mutation\_radius}
    \item \texttt{stagnation\_reset\_evals}
    \item \texttt{log\_path}
    \item \texttt{diagnostic\_chunk\_bytes}
    \item \texttt{report\_path} for detailed report output
\end{itemize}

For \texttt{"annealed\_hill\_climbing"}, \texttt{max\_mutation\_radius} is
normative and required.

Define baseline warmup count by
\[
w_{\mathrm{base}}:=\texttt{warmup\_baseline\_runs}\in\mathbb{N}_0.
\]
Warmup semantics are normative:
\begin{itemize}
    \item execute exactly $w_{\mathrm{base}}$ baseline warmup evaluations before
    the single normative baseline evaluation,
    \item warmup runs use the same evaluator path and timeout
    $\min(T_{\mathrm{eval}},T_{\mathrm{total}})$ as baseline evaluation,
    \item warmup runs are excluded from optimization metrics, incumbent updates,
    cache insertion, and the evaluation counter $k$,
    \item only the post-warmup baseline evaluation defines baseline
    $M(z_0^{\mathrm{can}})$ for optimization.
\end{itemize}
Warmup policy is evaluator-relevant and must be included in
$\eta_{\mathrm{eval}}$ and provenance.

Define stagnation-reset threshold by
\[
\kappa_{\mathrm{stag}}:=\texttt{stagnation\_reset\_evals}\in\mathbb{N}.
\]

Define the mutation-radius cap by
\[
r_{\mathrm{mut}}^{\max}:=\texttt{max\_mutation\_radius}\in\mathbb{N},
\qquad
r_{\mathrm{mut}}^{\max}\ge 1.
\]
This is an annealer execution-control parameter.

For the optional bounded same-task teacher-refresh wrapper around
\texttt{aiqi\_warmstart\_exact\allowbreak\_jh}, define
\[
R_{\mathrm{si}}:=\texttt{self\_improvement\_rounds}\in\mathbb{N},
\qquad
R_{\mathrm{si}}\ge 1,
\]
with default $R_{\mathrm{si}}=1$.
If $R_{\mathrm{si}}>1$, interpret $T_{\mathrm{total}}$ as the single
outer-loop wall-clock budget and use deterministic equal-split round deadlines
\[
\Delta_r^{\mathrm{si}}:=\frac{r}{R_{\mathrm{si}}}T_{\mathrm{total}},
\qquad r=0,\dots,R_{\mathrm{si}}.
\]
The warm-start exact-$\mathcal{J}_H$ controller then executes the bounded
same-task trace-refresh wrapper of
Section~\ref{sec:aiqi-warmstart-self-improve}; all rounds remain on the same
$(B,H,x)$ task and are evaluated against the same deployable objective
$\mathcal{J}_H$.

These controls affect runtime behavior and reproducibility, not the candidate search space.

\section{Bounds, Validity, and Canonicalization}

\subsection{Normative bounds}

In addition to backend/method inclusion and parameter ranges, the bounds model must include:
\begin{itemize}
    \item \texttt{max\_experts: usize},
    \item \texttt{max\_mixture\_nesting\_depth: usize}.
\end{itemize}

Strongly recommended optional bounds are:
\begin{itemize}
    \item \texttt{min\_experts}
    \item \texttt{allow\_duplicate\_experts}
    \item \texttt{required\_experts}
    \item \texttt{forbidden\_expert\_pairs}
\end{itemize}

\subsection{Candidate validity predicate}

Define $\mathrm{Valid}_{B,H}(z)=1$ iff all hold:
\begin{enumerate}
    \item schema conformance,
    \item bounds conformance,
    \item backend/library validation,
    \item referenced backend/method availability in the current build,
    \item every external reference in $z$ is a canonical identifier for a
    member of the shared side-information registry
    $\mathcal{S}_H^{\mathrm{side}}$; raw filesystem paths, URLs, and
    undeclared external assets are forbidden,
    \item mixture nesting depth $\le$ \texttt{max\_mixture\_nesting\_depth},
    \item per-node expert count in required interval.
\end{enumerate}

Invalid candidates are never evaluated.

\subsection{Canonical representation and cache key}\label{sec:canon-cache}

All candidates are canonicalized deterministically before hashing,
comparison, caching, evaluation, or output:
\begin{itemize}
    \item canonical backend naming,
    \item deterministic key ordering,
    \item order normalization only where semantics are order-insensitive,
    \item preservation of expert order where semantics are order-sensitive,
    \item removal of redundant defaults only when semantics-preserving,
    \item deterministic serialization bytes.
\end{itemize}

Let
\[
\chi_B:\mathcal{P}_{\mathrm{struct}}\to\{\mathrm{sensitive},\mathrm{insensitive}\}
\]
be the order-sensitivity classification over canonicalizable structure paths.
This classification is fixed by the backend schema/compiler contract for a
given build and bounds profile. Two compliant implementations must use the same
$\chi_B$; its identifier/version (or hash commitment) must be recorded in
provenance.

Define two distinct typed maps:
\[
\operatorname{canon}:\mathcal{Z}\to\mathcal{Z},
\qquad
\operatorname{ser}:\mathcal{Z}\to\{0,\dots,255\}^*.
\]
Here $\operatorname{canon}$ is an idempotent semantic canonicalizer
($\operatorname{canon}(\operatorname{canon}(z))=\operatorname{canon}(z)$), and
$\operatorname{ser}$ is deterministic byte serialization of a canonical
candidate.

Define canonical internal code map
\[
\kappa_B(z):=\operatorname{ser}(\operatorname{canon}(z)).
\]
This map is normative (not merely a cache utility): it defines
$L_B(z)=|\kappa_B(z)|$ in Section~\ref{sec:formal-problem} and objective
tie-breaking keys throughout this document. Required semantic properties are:
\begin{itemize}
    \item injective coding on canonical candidates
    $(\operatorname{canon}(z_1)=\operatorname{canon}(z_2)\iff
    \kappa_B(z_1)=\kappa_B(z_2))$,
    \item prefix-free code image over canonical candidates.
\end{itemize}

\begin{remark}[Side-information policy]
Let $H$ include a fixed shared side-information registry
\[
\mathcal{S}_H^{\mathrm{side}}=\{s_1,\dots,s_m\},
\]
available identically to encoder and decoder.

A candidate may reference shared side-information only through canonical
identifiers for members of $\mathcal{S}_H^{\mathrm{side}}$. The identifier, if
present, is part of the serialized candidate code $\kappa_B(z)$ and therefore
contributes to $L_B(z)$; the contents of the referenced shared asset do not
contribute to $L_B(z)$ because they are fixed by $H$.

Any candidate-specific external file, model weights, dictionary, URL, raw
filesystem path, or other undeclared external asset is forbidden. If evaluation
attempts to access such an asset, the candidate is invalid with diagnostic
reason \texttt{undeclared\_external\_asset}.
\end{remark}

Define canonical cache key
\[
\mathrm{Key}_{\mathrm{cache}}(z,\eta_{\mathrm{eval}},x):=
\Big(\operatorname{ser}(\operatorname{canon}(z)),\eta_{\mathrm{eval}},\operatorname{Hash}(x)\Big),
\]
where $\eta_{\mathrm{eval}}$ are evaluator-relevant execution controls and
$\operatorname{Hash}(x)$ is dataset identity hash. All cache lookups and inserts use this key.

Reference pseudocode:
\begin{verbatim}
function CanonAndSerialize(z):
    z_can <- canon(z)
    b     <- ser(z_can)
    return (z_can, b)

function CacheKey(z, eta_eval, x):
    (_, b) <- CanonAndSerialize(z)
    return (b, eta_eval, Hash(x))
\end{verbatim}

The deployable objective is operational.
The next section formalizes restricted mixture, MDL/MAP,
and optional universality semantics on the deployable semantic subclass.

\section{Restricted mixture semantics and truncated-universality dominance}
\label{sec:restricted-mixture}

This section formalizes three Solomonoff-adjacent statements for the deployable
candidate class:
\begin{enumerate}
    \item a restricted Bayesian/Solomonoff-style mixture over the bounded
    deployable semantic class,
    \item exact MDL/MAP equivalence and the dominant-term upper bound on bounded
    mixture codelength,
    \item an optional truncated-universality dominance theorem when the class
    contains a constant-overhead embedding of a truncated universal program
    family.
\end{enumerate}

\subsection{Semantic deployable class and code-length prior}

Operational and runtime lengths are byte-valued throughout this manuscript.
Implementation reports record these lengths in bytes. For MDL/MAP and mixture
statements in this semantic subsection, description lengths are expressed in
bits. Define the bit-normalized quantities
\[
L_B^{\mathrm{bit}}(z):=8L_B(z),
\qquad
C_x^{\mathrm{bit}}(z):=8C_x(z),
\qquad
\mathcal{J}_H^{\mathrm{bit}}(z;x):=L_B^{\mathrm{bit}}(z)+C_x^{\mathrm{bit}}(z).
\]
On this manuscript's byte-based operational convention,
$\mathcal{J}_H^{\mathrm{bit}}(z;x)=8\,\mathcal{J}_H(z;x)$ on
$\mathcal{F}^{\rho,\mathrm{dep}}(B,H)$.

\begin{remark}[Operational versus semantic normalization]
The operational objective $\mathcal{J}_H$ is specified in byte units. The
restricted-mixture and MDL/MAP statements use the bit-normalized objective
$\mathcal{J}_H^{\mathrm{bit}}$. On $\mathcal{F}^{\rho,\mathrm{dep}}(B,H)$,
$\mathcal{J}_H^{\mathrm{bit}}(z;x)=8\,\mathcal{J}_H(z;x)$, so minimizers are
identical.
\end{remark}

Let
\[
\mathcal{F}^{\rho,\mathrm{dep}}(B,H)
:=
\left\{
z\in\mathcal{F}^{\mathrm{dep}}(B,H):
\begin{aligned}[t]
&z \text{ induces a computable predictive semimeasure }\\
&\rho_z:\{0,\dots,255\}^*\to[0,1]
\end{aligned}
\right\}.
\]
Assume $\mathcal{F}^{\rho,\mathrm{dep}}(B,H)\neq\varnothing$.

For each $z\in\mathcal{F}^{\rho,\mathrm{dep}}(B,H)$, define the code-length
prior mass
\[
\pi_B(z):=2^{-L_B^{\mathrm{bit}}(z)}.
\]

Define the total semantic prior mass
\[
Z_B^{\rho,\mathrm{dep}}
:=
\sum_{z\in\mathcal{F}^{\rho,\mathrm{dep}}(B,H)} \pi_B(z).
\]

\begin{proposition}[Kraft bound on semantic deployable prior mass]
\label{prop:kraft-semantic}
If the exact internal code image is prefix-free, then
\[
0< Z_B^{\rho,\mathrm{dep}} \le 1.
\]
\end{proposition}

\begin{proof}
Positivity holds because
$\mathcal{F}^{\rho,\mathrm{dep}}(B,H)\neq\varnothing$ and each term
$2^{-L_B^{\mathrm{bit}}(z)}$ is strictly positive.
Because each byte is represented by a fixed 8-bit block, a prefix-free code
over byte strings induces a prefix-free binary code with bitlength
$L_B^{\mathrm{bit}}(z)=8L_B(z)$. The upper bound follows from the Kraft
inequality for prefix-free binary codes, applied to the semantic deployable
subclass.
\end{proof}

Define the \emph{unnormalized Solomonoff-style bounded mixture}
\[
\widetilde{\xi}_{B,H}(x)
:=
\sum_{z\in\mathcal{F}^{\rho,\mathrm{dep}}(B,H)}
2^{-L_B^{\mathrm{bit}}(z)}\rho_z(x).
\]

Define the \emph{normalized restricted Bayesian mixture}
\[
\xi_{B,H}(x)
:=
\sum_{z\in\mathcal{F}^{\rho,\mathrm{dep}}(B,H)}
w_B(z)\rho_z(x),
\qquad
w_B(z):=\frac{2^{-L_B^{\mathrm{bit}}(z)}}{Z_B^{\rho,\mathrm{dep}}}.
\]

\begin{proposition}[Semimeasure and mixture properties]
\label{prop:mixture-semimeasure}
If each $\rho_z$ is a semimeasure, then:
\begin{enumerate}
    \item $\widetilde{\xi}_{B,H}$ is a semimeasure,
    \item $\xi_{B,H}$ is a semimeasure,
    \item $\{w_B(z)\}_{z\in\mathcal{F}^{\rho,\mathrm{dep}}(B,H)}$ is a proper
    probability distribution on the semantic deployable class.
\end{enumerate}
\end{proposition}

\begin{proof}
Each claim follows from nonnegative linearity. Since
$Z_B^{\rho,\mathrm{dep}}\in(0,1]$ by
Proposition~\ref{prop:kraft-semantic},
the normalized weights are well-defined, nonnegative, and sum to $1$.
A nonnegative weighted sum of semimeasures is again a semimeasure.
\end{proof}

\begin{remark}[Scope of the restricted-mixture statement]
This outer restricted-mixture statement is a statement about the candidate class
as a whole, equipped with the exact code-length prior
$2^{-L_B^{\mathrm{bit}}(z)}$. It does
not follow merely from the fact that some individual candidates are themselves
Bayesian mixtures or ensembles internally.
\end{remark}

\subsection{MDL/MAP equivalence and dominant-term principle}

Assume in this subsection the idealized exact-coding semantics
\[
C_x^{\mathrm{bit}}(z) = -\log_2 \rho_z(x)
\qquad
\text{for each } z\in\mathcal{F}^{\rho,\mathrm{dep}}(B,H)
\]
Use extended-real convention
\[
-\log_2 0 := +\infty.
\]
Hence all MDL/MAP statements below are interpreted in
$\mathbb{R}_{\ge 0}\sqcup\{+\infty\}$.
Equivalently, the deployable objective on the semantic deployable class is
\[
\mathcal{J}_H^{\mathrm{bit}}(z;x)=L_B^{\mathrm{bit}}(z)-\log_2 \rho_z(x).
\]

\begin{proposition}[Exact MDL/MAP equivalence]
\label{prop:mdl-map}
For fixed data $x$, minimizing $\mathcal{J}_H^{\mathrm{bit}}(z;x)$ over
$z\in\mathcal{F}^{\rho,\mathrm{dep}}(B,H)$ is equivalent to maximizing
\[
2^{-L_B^{\mathrm{bit}}(z)}\rho_z(x),
\]
and equivalently to maximizing
\[
w_B(z)\rho_z(x).
\]
That is,
\[
\arg\min_{z\in\mathcal{F}^{\rho,\mathrm{dep}}(B,H)}\mathcal{J}_H^{\mathrm{bit}}(z;x)
=
\arg\max_{z\in\mathcal{F}^{\rho,\mathrm{dep}}(B,H)}
2^{-L_B^{\mathrm{bit}}(z)}\rho_z(x)
=
\arg\max_{z\in\mathcal{F}^{\rho,\mathrm{dep}}(B,H)}
w_B(z)\rho_z(x).
\]
\end{proposition}

\begin{proof}
For every $z$ in the semantic deployable class,
\[
\mathcal{J}_H^{\mathrm{bit}}(z;x)
=
L_B^{\mathrm{bit}}(z)-\log_2\rho_z(x)
=
-\log_2\!\bigl(2^{-L_B^{\mathrm{bit}}(z)}\rho_z(x)\bigr).
\]
For each $z$, one has
$2^{-L_B^{\mathrm{bit}}(z)}\rho_z(x)\in[0,1]$. With the convention
$-\log_2 0:=+\infty$, the map $y\mapsto -\log_2 y$ is strictly
order-reversing on $[0,1]$. Hence minimizing
$\mathcal{J}_H^{\mathrm{bit}}(z;x)$ is equivalent to maximizing
$2^{-L_B^{\mathrm{bit}}(z)}\rho_z(x)$.
Because $w_B(z)=2^{-L_B^{\mathrm{bit}}(z)}/Z_B^{\rho,\mathrm{dep}}$ and
$Z_B^{\rho,\mathrm{dep}}$ is constant in $z$, the same argmax is obtained for
$w_B(z)\rho_z(x)$.
\end{proof}

\begin{proposition}[Dominant-term upper bound for bounded mixture codelength]
\label{prop:dominant-term}
For every byte string $x$,
\[
-\log_2 \widetilde{\xi}_{B,H}(x)
\le
\inf_{z\in\mathcal{F}^{\rho,\mathrm{dep}}(B,H)} \mathcal{J}_H^{\mathrm{bit}}(z;x).
\]
Consequently,
\[
-\log_2 \xi_{B,H}(x)
\le
\inf_{z\in\mathcal{F}^{\rho,\mathrm{dep}}(B,H)} \mathcal{J}_H^{\mathrm{bit}}(z;x).
\]
\end{proposition}

\begin{proof}
For every fixed
$z\in\mathcal{F}^{\rho,\mathrm{dep}}(B,H)$,
\[
\widetilde{\xi}_{B,H}(x)
=
\sum_{z'}2^{-L_B^{\mathrm{bit}}(z')}\rho_{z'}(x)
\ge
2^{-L_B^{\mathrm{bit}}(z)}\rho_z(x)
=
2^{-\mathcal{J}_H^{\mathrm{bit}}(z;x)}.
\]
Applying $-\log_2$ yields
\[
-\log_2 \widetilde{\xi}_{B,H}(x)\le \mathcal{J}_H^{\mathrm{bit}}(z;x).
\]
Taking the infimum over $z$ proves the first inequality.

For the normalized mixture,
\[
\xi_{B,H}(x)=\frac{\widetilde{\xi}_{B,H}(x)}{Z_B^{\rho,\mathrm{dep}}},
\]
hence
\[
-\log_2 \xi_{B,H}(x)
=
-\log_2 \widetilde{\xi}_{B,H}(x)+\log_2 Z_B^{\rho,\mathrm{dep}}
\le
-\log_2 \widetilde{\xi}_{B,H}(x),
\]
because $Z_B^{\rho,\mathrm{dep}}\le 1$. Combine with the first inequality.
\end{proof}

\begin{remark}[Normalized-mixture strengthening]
Because $\xi_{B,H}(x)=\widetilde{\xi}_{B,H}(x)/Z_B^{\rho,\mathrm{dep}}$ with
$Z_B^{\rho,\mathrm{dep}}\le 1$, one has
$\xi_{B,H}(x)\ge\widetilde{\xi}_{B,H}(x)$ and therefore
$-\log_2\xi_{B,H}(x)\le-\log_2\widetilde{\xi}_{B,H}(x)$.
Thus the normalized mixture gives a bound at least as strong as the
unnormalized one.
\end{remark}

\begin{remark}[Meaning of the dominant-term statement]
Proposition~\ref{prop:dominant-term} gives the precise link between the exact
restricted mixture and the dominant-term MDL principle on the bounded class:
the best two-part code length upper bounds the bounded mixture codelength. No
reverse inequality with a constant additive gap is claimed without additional
assumptions.
\end{remark}

\subsection{Optional truncated-universality dominance}

Fix a universal prefix machine $U$ and truncation parameters
\[
K\in\mathbb{N},
\qquad
T_U>0,
\qquad
M_U>0.
\]

Let
\[
\mathcal{P}_{K,T_U,M_U}^H
\]
denote the set of all programs $p$ such that:
\begin{enumerate}
    \item $|p|\le K$,
    \item under the fixed profile $H$, program $p$ induces a computable
    predictive semimeasure $\nu_p^H$,
    \item the induced predictor is itself deployable under the same resource
    contract encoded by $H$ (including throughput and memory requirements),
    \item the program's universal-family execution remains within truncation
    limits $T_U$ and $M_U$.
\end{enumerate}

\begin{assumption}[Constant-overhead embedding of the truncated universal family]
\label{ass:universal-embed}
There exist a constant $c_U\ge 0$ and an injective map
\[
\iota:\mathcal{P}_{K,T_U,M_U}^H
\to
\mathcal{F}^{\rho,\mathrm{dep}}(B,H),
\qquad
p\mapsto z_p,
\]
such that for every $p\in\mathcal{P}_{K,T_U,M_U}^H$,
\[
\rho_{z_p}=\nu_p^H
\qquad\text{and}\qquad
L_B^{\mathrm{bit}}(z_p)\le |p|+c_U.
\]
\end{assumption}

\begin{theorem}[Truncated-universal dominance]
\label{thm:truncated-universal-dominance}
Under Assumption~\ref{ass:universal-embed}, for every byte string $x$,
\[
\widetilde{\xi}_{B,H}(x)
\ge
2^{-c_U}
\sum_{p\in\mathcal{P}_{K,T_U,M_U}^H}
2^{-|p|}\nu_p^H(x).
\]
Consequently,
\[
\xi_{B,H}(x)
\ge
2^{-c_U}
\sum_{p\in\mathcal{P}_{K,T_U,M_U}^H}
2^{-|p|}\nu_p^H(x).
\]
\end{theorem}

\begin{proof}
By definition of $\widetilde{\xi}_{B,H}$ and the inclusion
$z_p=\iota(p)\in\mathcal{F}^{\rho,\mathrm{dep}}(B,H)$,
\[
\widetilde{\xi}_{B,H}(x)
=
\sum_{z\in\mathcal{F}^{\rho,\mathrm{dep}}(B,H)}
2^{-L_B^{\mathrm{bit}}(z)}\rho_z(x)
\ge
\sum_{p\in\mathcal{P}_{K,T_U,M_U}^H}
2^{-L_B^{\mathrm{bit}}(z_p)}\rho_{z_p}(x).
\]
Using Assumption~\ref{ass:universal-embed},
\[
2^{-L_B^{\mathrm{bit}}(z_p)}\rho_{z_p}(x)
\ge
2^{-(|p|+c_U)}\nu_p^H(x)
=
2^{-c_U}2^{-|p|}\nu_p^H(x).
\]
Summing over $p$ proves the first inequality.

For the normalized mixture,
\[
\xi_{B,H}(x)=\widetilde{\xi}_{B,H}(x)/Z_B^{\rho,\mathrm{dep}}
\ge \widetilde{\xi}_{B,H}(x),
\]
because $Z_B^{\rho,\mathrm{dep}}\le 1$ by
Proposition~\ref{prop:kraft-semantic}. Hence the same lower bound holds for
$\xi_{B,H}$.
\end{proof}

\begin{corollary}[Truncated-universal deployable MDL baseline]
\label{cor:truncated-universal-baseline}
Under the exact-coding assumption of the previous subsection and
Assumption~\ref{ass:universal-embed}, for every byte string $x$,
\[
\inf_{z\in\mathcal{F}^{\rho,\mathrm{dep}}(B,H)}\mathcal{J}_H^{\mathrm{bit}}(z;x)
\le
c_U
+
\inf_{p\in\mathcal{P}_{K,T_U,M_U}^H}
\bigl(|p|-\log_2 \nu_p^H(x)\bigr).
\]
\end{corollary}

\begin{proof}
For each $p\in\mathcal{P}_{K,T_U,M_U}^H$, let $z_p=\iota(p)$. Then
\[
\mathcal{J}_H^{\mathrm{bit}}(z_p;x)
=
L_B^{\mathrm{bit}}(z_p)-\log_2\rho_{z_p}(x)
\le
|p|+c_U-\log_2 \nu_p^H(x).
\]
Taking the infimum over $p$ gives
\[
\inf_{z\in\mathcal{F}^{\rho,\mathrm{dep}}(B,H)}\mathcal{J}_H^{\mathrm{bit}}(z;x)
\le
\inf_{p\in\mathcal{P}_{K,T_U,M_U}^H}
\mathcal{J}_H^{\mathrm{bit}}(z_p;x)
\le
c_U
+
\inf_{p\in\mathcal{P}_{K,T_U,M_U}^H}
\bigl(|p|-\log_2 \nu_p^H(x)\bigr).
\]
\end{proof}

\begin{remark}[Meaning of the universality claim]
Theorem~\ref{thm:truncated-universal-dominance} is a class-richness statement.
It does not claim that the runtime optimizer must explicitly execute the
universal family on every run. It states that if the deployable candidate class
contains a constant-overhead encoding of the truncated universal family, then
the induced bounded mixture and the deployable MDL optimum inherit a formal
truncated-universality baseline.
\end{remark}

\section{Mutation Grammar and Search Neighborhood}

Let $\mathcal{N}_B(z)$ denote the legal finite neighborhood induced by bounds compiler $B$.
Mutations are generated from this grammar only, never by unconstrained raw JSON edits.

Allowed mutation families include: mixture-kind change, expert add/remove/replace, local parameter mutation, reorder, optional expert toggle, local nested mutation, incumbent substructure copy, and reset-to-incumbent.

For integer parameter $p\in[L,U]$, with elapsed fraction $u\in[0,1]$,
temperature $T(u)$, and mutation-radius cap $r_{\mathrm{mut}}^{\max}$, use local radius
\[
r(u)=\max\!\left(1,\left\lfloor r_{\mathrm{mut}}^{\max}\frac{T(u)}{T_0}\right\rfloor\right),
\]
then propose
\[
\Delta p \in \{-r(u),\ldots,-1,+1,\ldots,+r(u)\},
\]
clipped to $[L,U]$.

\subsection{Finite action compilation for planner exposure}\label{sec:action-compile}

For planner exposure, the bounded mutation grammar is compiled by a
deterministic compiler $\operatorname{compile}_B$ parameterized by $B$.
Define compiled action-cardinality
\[
N_B^{\mathrm{compile}}(B):=\big|\operatorname{compile}_B(B)\big|\in\mathbb{N},
\qquad
N_B^{\mathrm{compile}}(B)\ge 1,
\]
and planner-visible indexed action alphabet
\[
\mathcal{A}_B=\{0,\dots,N_B^{\mathrm{compile}}(B)-1\},
\qquad
|\mathcal{A}_B|=N_B^{\mathrm{compile}}(B).
\]
For planner-exposed controller configurations, cardinality consistency is
normative:
\[
N_B^{\mathrm{compile}}(B)=\texttt{agent\_actions}.
\]
After compiling $\mathcal{A}_B$ from $B$, initialization must verify
\[
|\mathcal{A}_B|=\texttt{agent\_actions}.
\]
On mismatch, initialization aborts with configuration error
\texttt{action\_alphabet\_mismatch}.
Each action index selects one compiled mutation template. The template decoder
is state-aware so that grammar families such as incumbent-copy,
reset-to-incumbent, and temperature-scaled local mutation remain representable.

Let $\mathcal{S}_B$ denote the annealer proposal-state space (including
proposal source, incumbent, elapsed time, and mutation RNG state), and let
$s_t\in\mathcal{S}_B$ denote the current proposal state. Action decoding is
specified by totalized state-aware semantics
\[
\mu_B^{\mathrm{tot}}:\mathcal{A}_B\times\mathcal{S}_B\to\mathcal{Z}\times\{0,1\},
\qquad
(a_t,s_t)\mapsto(\tilde z_t,\iota_t),
\]
where $\iota_t=1$ marks an inapplicable mutation primitive at state $s_t$.
If an implementation models decoder randomness explicitly rather than through
RNG state inside $s_t$, the equivalent form is a total Markov kernel
$\mu_B^{\mathrm{tot}}(\cdot\mid a_t,s_t)\in\Delta(\mathcal{Z}\times\{0,1\})$.

For this specification, the action-selection distribution is fixed to
\emph{uniform} over the compiled finite alphabet:
\[
\pi_T(a\mid s,B):=\frac{1}{N_B^{\mathrm{compile}}(B)},
\qquad
a\in\mathcal{A}_B.
\]
The subscript $T$ is retained for interface compatibility; in this
specification temperature dependence enters through state in
$\mu_B^{\mathrm{tot}}$ rather than through non-uniform action weights.

Define proposal-outcome kernel
\[
\widehat q_T(z,\iota\mid s,B)
:=
\sum_{a\in\mathcal{A}_B}
\pi_T(a\mid s,B)
\mathbf{1}\!\left\{\mu_B^{\mathrm{tot}}(a,s)=(z,\iota)\right\},
\]
and induced canonical-candidate marginal
\[
q_T(z'\mid s,B)
:=
\sum_{(z,\iota)\in\mathcal{Z}\times\{0,1\}}
\widehat q_T(z,\iota\mid s,B)
\mathbf{1}\!\left\{\operatorname{canon}(z)=z'\right\}.
\]
Here $s$ is the current annealer proposal state.
If stochastic decoder form is used, replace indicator terms above by the
corresponding kernel masses from $\mu_B^{\mathrm{tot}}(\cdot\mid a,s)$.

\section{Annealed Hill Climbing Algorithm}\label{sec:ahc-runtime}

\subsection{State and initialization}

State variables: current candidate $z_{\mathrm{cur}}$, best candidate
$z_{\mathrm{best}}$, elapsed wall time $t$, evaluation counter $k$, and
stagnation counter $n_{\mathrm{stag}}\in\mathbb{N}_0$.

Initialization requirements:
\begin{enumerate}
    \item load and canonicalize baseline candidate $z_0^{\mathrm{can}}:=\operatorname{canon}(z_0)$,
    \item execute $w_{\mathrm{base}}$ warmup baseline runs per the execution-control semantics above,
    \item evaluate $z_0^{\mathrm{can}}$ once normatively with timeout $\min(T_{\mathrm{eval}},T_{\mathrm{total}})$,
    \item enforce the baseline deployability precondition from
    Section~\ref{subsec:baseline-precondition}; abort immediately if baseline
    is non-deployable (fails success status, throughput floor, or memory cap),
    \item set $z_{\mathrm{cur}}=z_{\mathrm{best}}=z_0^{\mathrm{can}}$ and $n_{\mathrm{stag}}\leftarrow 0$.
\end{enumerate}

\subsection{Cooling schedule}

With elapsed fraction
\[
u=\min\!\left(1,\frac{t}{T_{\mathrm{total}}}\right),
\]
use log-linear cooling
\[
T(u)=T_{\min}\left(\frac{T_0}{T_{\min}}\right)^{1-u},
\]
default $T_0=1.0$, $T_{\min}=10^{-3}$.

\subsection{Transition rule}

At each loop iteration:
\begin{enumerate}
    \item compute remaining budget $B_t:=T_{\mathrm{total}}-t$;
    if $B_t\le 0$ then terminate immediately,
    \item compute $T(u)$,
    \item sample action $a_t\sim\pi_T(\cdot\mid s_t,B)$, decode
    $(\tilde z',\iota')=\mu_B^{\mathrm{tot}}(a_t,s_t)$, and canonicalize
    $\hat z':=\operatorname{canon}(\tilde z')$,
    \item if $\iota'=1$, assign $M(\hat z')$ status invalid with diagnostic
    reason \texttt{inapplicable\_action} without starting compression execution,
    \item else if $\hat z'$ fails $\mathrm{Valid}_{B,H}$, assign $M(\hat z')$
    status invalid with diagnostic reason \texttt{candidate\_invalid} without
    starting compression execution,
    \item else for cache lookup use
    $\mathrm{Key}_{\mathrm{cache}}(\hat z',\eta_{\mathrm{eval}},x)$;
    on cache miss evaluate with per-candidate timeout
    \[
    T_t^{\mathrm{eval}}:=\min(T_{\mathrm{eval}},B_t),
    \]
    \item set $K_{\mathrm{pre}}:=K(z_{\mathrm{best}})$; if
    $\mathcal{J}_H(\hat z';x)<+\infty$ and
    $K(\hat z')\prec_K K_{\mathrm{pre}}$, update
    $z_{\mathrm{best}}\leftarrow \hat z'$,
    \item if $\mathcal{J}_H(\hat z';x)<+\infty$, accept as current via Metropolis criterion:
    \[
    P_{\mathrm{accept}}=
    \begin{cases}
    1, & \mathcal{J}_H(\hat z';x)\le \mathcal{J}_H(z_{\mathrm{cur}};x),\\
    \exp\!\left(-\dfrac{\mathcal{J}_H(\hat z';x)-\mathcal{J}_H(z_{\mathrm{cur}};x)}{T(u)}\right), & \mathcal{J}_H(\hat z';x)>\mathcal{J}_H(z_{\mathrm{cur}};x),
    \end{cases}
    \]
    and non-deployable candidates ($\mathcal{J}_H(\hat z';x)=+\infty$) are rejected,
    setting $z_{\mathrm{cur}}\leftarrow\hat z'$ on acceptance,
    \item update stagnation counter by
    \[
    n_{\mathrm{stag}}\leftarrow
    \begin{cases}
    0,
    & \mathcal{J}_H(\hat z';x)<+\infty\ \land\ K(\hat z')\prec_K K_{\mathrm{pre}},\\
    n_{\mathrm{stag}}+1,
    & \mathcal{J}_H(\hat z';x)<+\infty\ \land\ \neg\big(K(\hat z')\prec_K K_{\mathrm{pre}}\big),\\
    n_{\mathrm{stag}},
    & \mathcal{J}_H(\hat z';x)=+\infty,
    \end{cases}
    \]
    \item if $n_{\mathrm{stag}}\ge\kappa_{\mathrm{stag}}$, set
    $z_{\mathrm{cur}}\leftarrow z_{\mathrm{best}}$ and
    $n_{\mathrm{stag}}\leftarrow 0$.
\end{enumerate}

\subsection{Termination}

Terminate when any condition holds:
\begin{itemize}
    \item elapsed wall-clock reaches the global deadline
    $T_{\mathrm{total}}$,
    \item \texttt{max\_evaluations} reached,
    \item neighborhood exhaustion is proven,
    \item unrecoverable evaluator failure occurs.
\end{itemize}

Hard-budget enforcement is normative: no candidate evaluation may run with a
timeout exceeding remaining global budget, and no new blocking evaluation may
start once $B_t\le 0$.

Output must include:
\begin{itemize}
    \item deterministic serialized $z_{\mathrm{best}}$ at \texttt{output\_config\_path},
    \item structured report at optional \texttt{report\_path}.
\end{itemize}

\section{Reward and Observation (Future Environment Compatibility)}\label{sec:reward-observation}

This section specifies the single normative planner-reward path.
Candidate validity, evaluator semantics, objective-feasibility semantics
(including throughput and memory deployability constraints), and output
definition remain shared across controller choices.

\subsection{Objective-aligned planner reward path}

For exact objective alignment under planner controllers, use
incumbent-objective-difference reward
\[
R_t^{\mathcal{J}}:=\mathcal{J}_H(b_t;x)-\mathcal{J}_H(b_{t+1};x),
\]
where $\mathcal{J}_H$ is the deployable objective from
Section~\ref{sec:formal-problem}. This preserves exact objective ordering.

For fixed finite horizon $m\in\mathbb{N}$, the exact cumulative-improvement
target is
\[
G_t^{(m)}:=\sum_{k=0}^{m-1}R_{t+k}^{\mathcal{J}}
=\mathcal{J}_H(b_t;x)-\mathcal{J}_H(b_{t+m};x),
\]
so finite-horizon planner targets are exact objective decreases over $m$
executed steps.

With incumbent $b$ before evaluating $z$, define diagnostic deltas
\[
\delta_c(z\mid b)=\frac{c(b)-c(z)}{\max(c(b),1)},
\qquad
\delta_v(z\mid b)=\frac{\tau(b)-\tau(z)}{\max(\tau(b),\varepsilon)},
\]
with fixed $\varepsilon>0$ (default $\varepsilon=10^{-9}$).

For total observation semantics, define
\[
n_+:=\max(|x|,1).
\]
Use reserved failure symbols
\[
\bot_c,\bot_\ell,\bot_\tau,\bot_{\delta_c},\bot_{\delta_v}.
\]

Let $\Sigma_B$ be a finite signature alphabet and define deterministic
signature map
\[
\sigma:\mathcal{Z}\to\Sigma_B.
\]

Define totalized diagnostics:
\[
\tilde c(z)=
\begin{cases}
 c(z)/n_+, & \mathrm{status}(z)=\mathrm{success},\\
 \bot_c, & \text{otherwise,}
\end{cases}
\qquad
\tilde \ell(z)=
\begin{cases}
 \ell(z)/n_+, & \mathrm{status}(z)=\mathrm{success},\\
 \bot_\ell, & \text{otherwise,}
\end{cases}
\]
\[
\tilde \tau(z)=
\begin{cases}
\tau(z)/T_{\mathrm{eval}}, & \mathrm{status}(z)=\mathrm{success},\\
1, & \mathrm{status}(z)=\mathrm{timeout},\\
\bot_\tau, & \text{otherwise,}
\end{cases}
\]
\[
\tilde\delta_c(z\mid b)=
\begin{cases}
\delta_c(z\mid b), & \mathrm{status}(z)=\mathrm{success},\\
\bot_{\delta_c}, & \text{otherwise,}
\end{cases}
\qquad
\tilde\delta_v(z\mid b)=
\begin{cases}
\delta_v(z\mid b), & \mathrm{status}(z)=\mathrm{success},\\
\bot_{\delta_v}, & \text{otherwise.}
\end{cases}
\]

Define typed value spaces
\[
\mathcal{V}_c:=\mathbb{R}_{\ge 0}\sqcup\{\bot_c\},
\quad
\mathcal{V}_{\ell}:=\mathbb{R}_{\ge 0}\sqcup\{\bot_\ell\},
\quad
\mathcal{V}_{\tau}:=[0,1]\sqcup\{\bot_\tau\},
\]
\[
\mathcal{V}_{\delta_c}:=\mathbb{R}\sqcup\{\bot_{\delta_c}\},
\qquad
\mathcal{V}_{\delta_v}:=\mathbb{R}\sqcup\{\bot_{\delta_v}\}.
\]
Then raw observation codomain is
\[
\mathcal{O}^{\mathrm{raw}}_B:=
\{0,1\}\times\mathcal{V}_c\times\mathcal{V}_{\ell}\times\mathcal{V}_{\tau}
\times\mathcal{V}_{\delta_c}\times\mathcal{V}_{\delta_v}\times\Sigma_B\times\{0,1\}.
\]

The raw observation payload is
\[
o^{\mathrm{raw}}(z\mid b,d_t)=
\Big(
\mathbf{1}_{\mathrm{fail}},
\tilde c(z),
\tilde \ell(z),
\tilde \tau(z),
\tilde\delta_c(z\mid b),
\tilde\delta_v(z\mid b),
\sigma(z),
d_t
\Big),
\qquad
o^{\mathrm{raw}}(z\mid b,d_t)\in\mathcal{O}^{\mathrm{raw}}_B,
\]
where $\mathbf{1}_{\mathrm{fail}}=1$ for timeout/invalid/error or for successful candidates with
$\theta(z)<\theta_{\min}$ or $\mu(z)>\mu_{\max}$, and $0$ otherwise,
and $d_t\in\{0,1\}$ is terminal indicator emitted by this environment.

Using planner parameters
$(L_O,b_O,\mathcal{R}_B,r_{\min},r_{\max})$ from
Section~\ref{sec:aixi-embedding}, define deterministic adapter maps:
\[
\Pi_R:\mathbb{R}\to\mathcal{R}_B,
\qquad
\Pi_R(r):=\operatorname{clip}_{[r_{\min},r_{\max}]}
\!\left(\left\lfloor r+\tfrac{1}{2}\right\rfloor\right),
\]
\[
\Psi_O:=\Psi_O^{\eta_O},
\qquad
\Psi_O^{\eta_O}:\mathcal{O}^{\mathrm{raw}}_B\to\mathcal{O}_B^{L_O},
\]
for a fixed observation-adapter specification
\[
\eta_O:=(\text{field order},\text{sentinel coding},\text{quantization/clipping},\text{packing rule},L_O,b_O).
\]
where
\[
\operatorname{clip}_{[L,U]}(x):=\min\!\big(\max(x,L),U\big).
\]
Admissibility requirements are
\[
\Pi_R(r)\in [r_{\min},r_{\max}]\cap\mathbb{Z}
\quad\text{for all }r\in\mathbb{R},
\]
\[
0\le (\Psi_O(o))_i\le 2^{b_O}-1
\quad\text{for all }o\in\mathcal{O}^{\mathrm{raw}}_B,\ i\in\{1,\dots,L_O\},
\]
with the degenerate case $(\Psi_O(o))_i=0$ when $b_O=0$.
For reproducible planner behavior, $\eta_O$ is fixed for a run and must be
recorded (or hash-committed) in provenance.

\begin{remark}[Adapter boundary]
The maps $\Pi_R$ and $\Psi_O$ are the generic raw-environment interfaces in
this section. The exact planner-state embedding in
Section~\ref{sec:aixi-embedding} uses the stronger specialized notation
$\Psi_H$ and $Q_R$: there the observation channel encodes planner state
directly and the reward already lies in $\mathcal{R}_B$. This is a
specialization of the same objective-difference reward path, not a change in
planner semantics.
\end{remark}

\section{Core Properties}

\begin{proposition}[Feasibility safety]
Any returned $z_{\mathrm{best}}$ is a deployable evaluated candidate satisfying
bounds, timeout feasibility, minimum-throughput, and memory-cap constraints.
\end{proposition}

\begin{proof}
$z_{\mathrm{best}}$ starts at a deployable baseline and is updated only with
deployable candidates under $K$. Timeout/invalid/error and successful but
non-deployable candidates never enter $\mathcal{S}_t$.
\end{proof}

\begin{proposition}[Monotone incumbent key]
Across best updates, the incumbent key $K(z_{\mathrm{best}})$ is strictly decreasing in lexicographic order.
\end{proposition}

\begin{proof}
By update rule, $z_{\mathrm{best}}$ changes only when $K(z')\prec_K K(z_{\mathrm{best}})$. Hence each update is a strict decrease in $K$.
\end{proof}

\begin{proposition}[Objective monotonicity under key updates]
Across best updates, incumbent objective is non-increasing: if
$z_{\mathrm{best}}$ changes from $z$ to $z'$, then
$\mathcal{J}_H(z';x)\le\mathcal{J}_H(z;x)$.
\end{proposition}

\begin{proof}
The incumbent updates only when
$K(z')=(\mathcal{J}_H(z';x),\operatorname{ser}(\operatorname{canon}(z')))$ is
lexicographically smaller than
$K(z)=(\mathcal{J}_H(z;x),\operatorname{ser}(\operatorname{canon}(z)))$.
Hence $\mathcal{J}_H(z';x)\le\mathcal{J}_H(z;x)$.
\end{proof}

\begin{remark}[Conditional replay determinism guarantee]
If evaluator behavior, timeout clock semantics, RNG stream, build/features, and
execution controls are fixed and deterministic, then the tuning trace and final
output are deterministic.
\end{remark}

\begin{remark}
Wall-clock jitter can perturb timeout boundaries in practice. Therefore reproducibility is a normative requirement and should be recorded with timing and environment metadata.
\end{remark}

\section{Annealed-Controller Assumptions and Non-Asymptotic Guarantees}\label{sec:annealed-guarantees}

Unless explicitly stated otherwise, this section applies to the
\texttt{"annealed\_hill\_climbing"} runtime path from
Section~\ref{sec:ahc-runtime}.

\begin{assumption}[Finite grammar and computable validity]\label{ass:finite-grammar-validity}
For every valid candidate $z$, the neighborhood $\mathcal{N}_B(z)$ is finite and effectively enumerable; canonicalization and validity checks terminate.
\end{assumption}

\begin{assumption}[Positive proposal support for annealed controller]\label{ass:positive-proposal-support}
Let $z_{\mathrm{src}}(s)$ be the proposal-source component of state $s$. For
each reachable proposal state $s$ and each
$z'\in \mathcal{N}_B\big(z_{\mathrm{src}}(s)\big)$, there exists at least one action
$a\in\mathcal{A}_B$ and candidate $z\in\mathcal{Z}$ such that either
\[
\mu_B^{\mathrm{tot}}(a,s)=(z,0)
\]
in deterministic form, or
\[
\mu_B^{\mathrm{tot}}\big((z,0)\mid a,s\big)>0
\]
in stochastic form, and $\operatorname{canon}(z)=z'$. Under the specified
uniform policy
$\pi_T(a\mid s,B)=1/N_B^{\mathrm{compile}}(B)$, this implies
$q_T(z'\mid s,B)>0$.
\end{assumption}

\begin{assumption}[Per-iteration overhead lower bound]\label{ass:iter-overhead-lb}
There exists $h_{\min}>0$ such that each loop iteration consumes at least $h_{\min}$ wall-clock seconds including proposal, canonicalization, validation, cache lookup, and bookkeeping.
\end{assumption}

\begin{proposition}[Finite-horizon iteration bound]\label{prop:finite-iter-bound}
Let $N_{\mathrm{iter}}$ be total loop iterations before termination by time budget $T_{\mathrm{total}}$. Under Assumption~\ref{ass:iter-overhead-lb},
\[
N_{\mathrm{iter}}\le\left\lfloor\frac{T_{\mathrm{total}}}{h_{\min}}\right\rfloor.
\]
Hence all diagnostic counts (attempted proposals, invalids, cache hits, cache misses, successes) are finite and budget-bounded.
\end{proposition}

\begin{proposition}[Annealed finite-horizon improvement probability lower bound]\label{prop:finite-horizon-improve}
Fix horizon $H\in\mathbb{N}$. Suppose at each of the first $H$ deployable
evaluation opportunities, conditional on past history, the probability of
producing a strictly better incumbent is at least $p_*>0$. Then
\[
\mathbb{P}(\text{at least one incumbent improvement by }H)\ge 1-(1-p_*)^H.
\]
Consequently, the probability of no improvement by horizon $H$ decays at least geometrically.
\end{proposition}

\begin{proposition}[Non-asymptotic acceptance envelope]\label{prop:accept-envelope}
For a deployable uphill proposal with objective gap
$\Delta\mathcal{J}=\mathcal{J}_H(z';x)-\mathcal{J}_H(z_{\mathrm{cur}};x)>0$ at
elapsed fraction $u$, acceptance probability is
\[
\mathbb{P}_{\mathrm{accept}}(u)=\exp\!\left(-\frac{\Delta\mathcal{J}}{T(u)}\right).
\]
With log-linear cooling and $T_{\min}\le T(u)\le T_0$, this implies
\[
\exp\!\left(-\frac{\Delta\mathcal{J}}{T_{\min}}\right)
\le
\mathbb{P}_{\mathrm{accept}}(u)
\le
\exp\!\left(-\frac{\Delta\mathcal{J}}{T_0}\right),
\]
and $\mathbb{P}_{\mathrm{accept}}(u)$ is monotone decreasing in elapsed fraction
$u$ (strictly on $u\in[0,1)$).
\end{proposition}

\begin{remark}
These are finite-time statements: they do not claim global optimality under finite budgets. They characterize the expected behavior envelope and make engineering trade-offs explicit.
\end{remark}

\subsection{Idealized asymptotic theorem (explicitly out-of-budget scope)}\label{sec:annealed-asymptotic-scope}

The guarantees above are finite-time and budget-respecting. The theorem below
is deliberately scope-delimited: it describes an idealized infinite-opportunity
annealed process and is not a production guarantee under finite-budget runtime.

\begin{assumption}[Finite feasible candidate key space]\label{ass:finite-success-key-space}
The deployable set $\mathcal{F}^{\mathrm{dep}}(B,H)$ is finite and
nonempty.
\end{assumption}

\begin{assumption}[Infinite deployable opportunities]\label{ass:infinite-success-opportunities}
Index deployable evaluation opportunities by $h\in\mathbb{N}$ and let $b_h$
denote the incumbent after opportunity $h$. The process admits infinitely many
such opportunities almost surely.
\end{assumption}

\begin{assumption}[Uniform one-step improvement support]\label{ass:uniform-improvement-support}
There exists $p_{\mathrm{imp}}\in(0,1]$ such that for every deployable opportunity
index $h$, conditional on any realized history with incumbent $b_h$ not globally
minimal in $K$ over $\mathcal{F}^{\mathrm{dep}}(B,H)$, the probability of a strict
incumbent-key improvement at the next deployable opportunity is at least
$p_{\mathrm{imp}}$.
\end{assumption}

Define incumbent rank by
\[
\operatorname{rank}_K(b):=
\big|\{z\in\mathcal{F}^{\mathrm{dep}}(B,H):K(z)\prec_K K(b)\}\big|.
\]
Then $\operatorname{rank}_K(b)=0$ iff $b$ is globally $K$-minimal.

\begin{theorem}[Idealized almost-sure eventual global-optimum hit]\label{thm:idealized-asymptotic-hit}
Under Assumptions~\ref{ass:finite-success-key-space}--\ref{ass:uniform-improvement-support},
\[
\mathbb{P}\!\big(\exists h\in\mathbb{N}:\operatorname{rank}_K(b_h)=0\big)=1.
\]
Moreover, for any realized history with $\operatorname{rank}_K(b_h)>0$ and any
$H\in\mathbb{N}$,
\[
\mathbb{P}\!\big(\operatorname{rank}_K(b_{h+j})=\operatorname{rank}_K(b_h)\ \forall j=1,\dots,H\mid\mathcal{H}_h\big)
\le (1-p_{\mathrm{imp}})^H,
\]
where $\mathcal{H}_h$ denotes history up to deployable opportunity $h$.
\end{theorem}

\begin{proof}
If $\operatorname{rank}_K(b_h)>0$, Assumption~\ref{ass:uniform-improvement-support}
implies conditional probability at least $p_{\mathrm{imp}}$ of a strict rank
drop at the next deployable opportunity. Therefore the probability of no rank
drop for $H$ consecutive deployable opportunities is at most
$(1-p_{\mathrm{imp}})^H$, proving the finite-horizon bound. Letting
$H\to\infty$ gives probability $0$ of never dropping rank from any positive
rank state. Since rank is a nonnegative integer that strictly decreases on
every incumbent improvement and is bounded above by
$|\mathcal{F}^{\mathrm{dep}}(B,H)|-1$, only finitely many drops are needed to
reach rank $0$. Hence rank $0$ is reached almost surely.
\end{proof}

\begin{remark}
Theorem~\ref{thm:idealized-asymptotic-hit} is an idealized asymptotic statement.
Finite-budget runtime with hard deadline $T_{\mathrm{total}}$ and optional
\texttt{max\_evaluations} remains governed by the non-asymptotic statements in
this section.
\end{remark}

\section{Computational Complexity}\label{sec:computational-complexity}

Let:
\begin{itemize}
    \item $n=|x|$ dataset bytes,
    \item $N$ attempted proposals,
    \item $U\le N$ unique canonical candidates evaluated or cache-hit,
    \item $|z|$ canonical candidate representation size,
    \item $C_{\mathrm{comp}}(n,z)$ compression runtime.
\end{itemize}

Per proposal attempt, overhead is
\[
O\big(C_{\mathrm{mut}}(z,B)+C_{\mathrm{canon}}(|z|)+C_{\mathrm{valid}}(z,B)+C_{\mathrm{hash}}(|z|)\big),
\]
with $C_{\mathrm{canon}}(|z|)=O(|z|\log|z|)$ under key sorting, and $C_{\mathrm{hash}}(|z|)=O(|z|)$.

For cache misses, evaluation cost is bounded by hard timeout:
\[
C_{\mathrm{eval}}(z) \le
\min\{C_{\mathrm{comp}}(n,z),T_t^{\mathrm{eval}}\}+C_{\mathrm{kill}},
\qquad
T_t^{\mathrm{eval}}\le T_{\mathrm{total}}-t.
\]

Hence evaluation runtime is globally deadline-bounded by design, with only
constant finalization overhead after the deadline check:
\[
\mathrm{Runtime} \le T_{\mathrm{total}} + O(C_{\mathrm{finalize}}).
\]

Memory complexity is dominated by cache:
\[
O\big(U\cdot(|\operatorname{ser}(\operatorname{canon}(z))|+|M(z)|)\big).
\]

For the MC-AIXI/AIQI-family runtime path, let $N_{\mathrm{dec}}$ be the number of
real decision steps and let
\[
N_{\mathrm{sim}}:=\sum_{t=0}^{N_{\mathrm{dec}}-1}N_{\mathrm{sim},t}
\]
be total planner-internal simulations, where $N_{\mathrm{sim},t}$ is
simulations at decision step $t$ (bounded by
\texttt{planner\_simulations\_per\_step}). If
$C_{\mathrm{sim}}(t)$ denotes per-simulation planner compute at step $t$, then
controller-internal planning overhead is
\[
O\!\left(\sum_{t=0}^{N_{\mathrm{dec}}-1}N_{\mathrm{sim},t}\,C_{\mathrm{sim}}(t)\right)
=O\big(N_{\mathrm{sim}}\,\overline C_{\mathrm{sim}}\big),
\]
with $\overline C_{\mathrm{sim}}$ the average per-simulation planner cost.
This term affects search-control compute only and does not alter evaluator,
validity, or output semantics.

\section{Implementation Constraints for \texttt{infotheory}}

\subsection{Semantic equivalence to normal compression path}

Evaluation must preserve the same semantics as normal compression invocation (e.g., the existing mixture JSON path via \texttt{--rate-backend mixture} with \texttt{--method} file input). Library-internal calls are acceptable only if behavior is equivalent.

\subsection{CLI exposure}

Expose capability as
\begin{verbatim}
infotheory tune <spec.json>
\end{verbatim}
with file-based specification, consistent with other JSON-configured modes.

\subsection{Required report and provenance fields}

The structured report should record at minimum:
\begin{itemize}
    \item baseline and best metrics ($c,\tau,\mu,\mathcal{J}_H$),
    \item operational length quantities are recorded in bytes
    ($c,L_B,C_x,\mathcal{J}_H$),
    \item any optional semantic bit-normalized quantities are explicitly
    labeled and computed as $8\times$ the corresponding byte quantities,
    \item deltas and evaluation counters,
    \item timeout, invalid, and successful-but-nondeployable counts,
    \item final best objective,
    \item final best-move reward,
    \item crate/tool version,
    \item feature set (if available),
    \item seed,
    \item configured throughput floor $\theta_{\min}$ and derived runtime cap $\tau_{\theta}=n/\theta_{\min}$,
    \item configured memory cap $\mu_{\max}$,
    \item dataset path and dataset content hash,
    \item baseline config hash,
    \item initial teacher-dataset provenance/hash (if used),
    \item warmup policy and warmup count $w_{\mathrm{base}}$,
    \item self-improvement round count $R_{\mathrm{si}}$ and realized per-round trace counts (if enabled),
    \item bounds hash,
    \item canonicalization order-classification identifier/version (or hash commitment),
    \item observation-adapter or exact-state encoder specification (or hash commitment),
    \item output config hash.
\end{itemize}

\subsection{Cache key requirements}

Use the canonical key definition from Section~\ref{sec:canon-cache}:
\[
\mathrm{Key}_{\mathrm{cache}}(z,\eta_{\mathrm{eval}},x)=
\Big(\operatorname{ser}(\operatorname{canon}(z)),\eta_{\mathrm{eval}},\operatorname{Hash}(x)\Big).
\]

For all controller paths, before evaluating a canonical candidate $\hat z$, the runtime must query the cache key above instantiated at $(\hat z,\eta_{\mathrm{eval}},x)$. On cache hit, the cached evaluator output $M(\hat z)$ is reused and no new compression execution occurs. Any state-relative reward or observation components are recomputed from that cached evaluator output and the current runtime state. This ensures reasonable performance, without affecting behaviour under hardware determinism assumptions.

\section{Expected Theoretical Behavior}\label{sec:expected-theoretical-behavior}

Under fixed hard budgets and bounded mutation grammar, this method is a practical anytime algorithm:
\begin{itemize}
    \item it is guaranteed to return a deployable incumbent (baseline or better),
    \item it enforces throughput and memory limits as hard deployability constraints,
    \item it never promotes infeasible outcomes,
    \item it trades early exploration for late exploitation through temperature decay,
    \item it preserves exact objective semantics via $\mathcal{J}_H$ on the deployable region.
\end{itemize}

The exact-$\mathcal{J}_H$ statement applies to the annealed hill-climbing
controller and to planner controllers under the objective-aligned
objective-difference reward path from
Section~\ref{sec:reward-observation}.

No claim of global optimality is made under finite budgets. The exact guarantee is that uphill moves have strictly positive acceptance probability under annealing:
\[
\Delta\mathcal{J}>0,\ T(u)>0
\implies
\mathbb{P}_{\mathrm{accept}}(u)=\exp(-\Delta\mathcal{J}/T(u))>0,
\]
whereas deterministic hill climbing assigns probability $0$ to uphill moves.

The asymptotic claims in this document are:
\begin{itemize}
    \item the annealed idealized eventual-hit theorem
    (Theorem~\ref{thm:idealized-asymptotic-hit}), conditional on
    Section~\ref{sec:annealed-asymptotic-scope}, and
    \item the planner-side MC-AIXI asymptotic convergence result
    (Theorem~\ref{thm:planner-convergence-mcaixi}, with
    Proposition~\ref{prop:planner-convergence-optset} giving the broader
    set-valued non-unique optimal-sequence variant), conditional on
    Section~\ref{sec:planner-convergence-mcaixi}.
\end{itemize}
Both are assumption-conditional asymptotic statements and are outside
finite-budget runtime guarantees.

\section{Exact planner-state embedding for MC-AIXI(FAC-CTW) under a fixed hardware profile}\label{sec:aixi-embedding}

This section formalizes a planner-facing environment whose dynamics are Markov
by construction, while separating two logically distinct use-cases:
\begin{enumerate}
    \item a budget-free finite-horizon planning variant for direct
    $\rho$UCT / MC-AIXI consistency statements, and
    \item an optional reset-augmented continuing variant for stationary/ergodic
    finite-state Markov analysis.
\end{enumerate}
All finite-MDP and observed-Markov claims in this section are conditional on
Assumptions~\ref{ass:finite-Z}--\ref{ass:no-hidden}.

\subsection{Fixed-profile semantics}

Fix a hardware/evaluator profile
\[
\begin{aligned}
H=\bigl(&\text{machine class},\text{CPU pinning policy},\text{RAM cap},\\
&\text{build/features},\eta_{\mathrm{eval}},\\
&\text{wall-clock timing methodology},\text{seed policy}\bigr).
\end{aligned}
\]
All planner-facing semantics below are defined relative to $H$.

\begin{assumption}[Finite planner-exposed candidate domain]\label{ass:finite-Z}
The planner-exposed canonical candidate set
\[
\mathcal{Z}_H^{\mathrm{can}} \subseteq \mathcal{Z}
\]
is finite.
\end{assumption}

\begin{remark}
Assumption~\ref{ass:finite-Z} is normative for this embedding. Any parameter
that would otherwise range over an infinite or continuous set must be
explicitly discretized before planner exposure; this discretization is part of
$B$ and part of the canonicalization contract.
\end{remark}

Define status alphabet
\[
\mathcal{S}_{\mathrm{status}}
:=
\{\mathrm{success},\mathrm{timeout},\mathrm{invalid},\mathrm{error}\}.
\]
Let $\mathcal{D}_H$ be a finite diagnostic codomain containing distinguished
token
\[
d_{\mathrm{inapp}}:=\texttt{inapplicable\_action},
\]
and let $\bot_c,\bot_\ell,\bot_\tau,\bot_\mu$ denote failure sentinels for
undefined size/log-loss/time/memory fields. Define evaluator codomain
\[
\mathcal{Y}_H
:=
\mathcal{S}_{\mathrm{status}}
\times (\mathbb{N}\sqcup\{\bot_c\})
\times (\mathbb{R}_{\ge 0}\sqcup\{\bot_\ell\})
\times (\mathbb{R}_{\ge 0}\sqcup\{\bot_\tau\})
\times (\mathbb{R}_{\ge 0}\sqcup\{\bot_\mu\})
\times \mathcal{D}_H.
\]

\begin{assumption}[Profile-fixed semantic evaluator]\label{ass:det-eval}
There exists a deterministic semantic evaluator
\[
M_H^\circ:\mathcal{Z}_H^{\mathrm{can}}\to\mathcal{Y}_H,
\qquad
M_H^\circ(z)=\big(\mathrm{status}_H(z),c_H(z),\ell_H(z),\tau_H(z),\mu_H(z),\operatorname{diag}_H(z)\big),
\]
with:
\begin{itemize}
    \item $\mathrm{status}_H(z)\in\mathcal{S}_{\mathrm{status}}$,
    \item $c_H(z)\in\mathbb{N}$ on success and $c_H(z)=\bot_c$ otherwise,
    \item $\ell_H(z)\in\mathbb{R}_{\ge 0}$ on success and
    $\ell_H(z)=\bot_\ell$ otherwise,
    \item $\tau_H(z)\in\mathbb{R}_{\ge 0}$ on success and
    $\tau_H(z)=\bot_\tau$ otherwise,
    \item $\mu_H(z)\in\mathbb{R}_{\ge 0}$ on success and
    $\mu_H(z)=\bot_\mu$ otherwise,
    \item $\operatorname{diag}_H(z)\in\mathcal{D}_H$.
\end{itemize}
This is the fixed-profile specialization of the six-field evaluator contract
from Section~\ref{sec:formal-problem}. Here $\tau_H(z)$ denotes elapsed
wall-clock time in the same unit family as $\tau(z)$, $T_{\mathrm{eval}}$,
and $T_{\mathrm{total}}$. Any additional deterministic profile costs (for
example, PMU cycles) may be recorded in $\operatorname{diag}_H(z)$, but are
not denoted by $\tau_H(z)$.
\end{assumption}

For successful candidates under profile $H$, define throughput
\[
\theta_H(z):=
\begin{cases}
\dfrac{n}{\tau_H(z)}, & \mathrm{status}_H(z)=\mathrm{success},\ \tau_H(z)>0,\\[4pt]
+\infty, & \mathrm{status}_H(z)=\mathrm{success},\ \tau_H(z)=0.
\end{cases}
\]

\begin{remark}
Assumption~\ref{ass:det-eval} does not remove hardware from the model. It
incorporates hardware methodology into $H$; once $H$ is fixed, evaluator
semantics are required to be deterministic relative to that profile.
\end{remark}

\begin{assumption}[Deployable baseline initialization]\label{ass:baseline-succ}
This is the fixed-profile restatement of the normative input precondition from
Section~\ref{subsec:baseline-precondition}.

In a conforming runtime, $b_0$ is the incumbent produced by the single
post-warmup evaluation of $z_0^{\mathrm{can}}$.

The initial incumbent $b_0\in\mathcal{Z}_H^{\mathrm{can}}$ satisfies
\[
\mathrm{status}_H(b_0)=\mathrm{success}
\quad\land\quad
\theta_H(b_0)\ge\theta_{\min}
\quad\land\quad
\mu_H(b_0)\le\mu_{\max}.
\]
\end{assumption}

Define the deployable canonical subset
\[
\mathcal{Z}_{H,\mathrm{dep}}^{\mathrm{can}}:=
\left\{z\in\mathcal{Z}_H^{\mathrm{can}}:\ \mathrm{status}_H(z)=\mathrm{success}\ \land\ \theta_H(z)\ge\theta_{\min}\ \land\ \mu_H(z)\le\mu_{\max}\right\}.
\]
By Assumption~\ref{ass:baseline-succ}, this set is nonempty.

Define profile-fixed two-part objective components
\[
C_x(z):=c_H(z),
\qquad
L_B(z):=\big|\operatorname{ser}(z)\big|,
\]
and objective
\[
\mathcal{J}_H(z;x):=
\begin{cases}
L_B(z)+C_x(z),
& z\in\mathcal{Z}_{H,\mathrm{dep}}^{\mathrm{can}},\\[3pt]
+\infty,
& \text{otherwise.}
\end{cases}
\]
This enforces a single deployability-gated objective in the planner path:
non-deployable candidates receive $+\infty$ exactly as timeout/invalid/error
outcomes.

Define incumbent key
\[
K_H(z):=\big(\mathcal{J}_H(z;x),\operatorname{ser}(z)\big),
\]
ordered lexicographically.

\subsection{State, action, and omitted-variable discipline}

Fix planner-interface integer parameters
\[
\begin{aligned}
N_B&:=\texttt{agent\_actions}=N_B^{\mathrm{compile}}(B)\in\mathbb{N},\ N_B\ge 1,\\
b_O&:=\texttt{observation\_bits}\in\mathbb{N}_0,\\
L_O&:=\max(1,\texttt{observation\_stream\_len}),\\
b_R&:=\texttt{reward\_bits}\in\mathbb{N},\ b_R\ge 1,\\
r_{\min}&:=\texttt{min\_reward}\in\mathbb{Z},\\
r_{\max}&:=\texttt{max\_reward}\in\mathbb{Z},\ r_{\max}>r_{\min},\\
\rho_{\mathrm{off}}&:=\texttt{reward\_offset}\in\mathbb{Z}.
\end{aligned}
\]
Define
\[
\mathcal{A}_B=\{0,\dots,N_B-1\},
\qquad
\mathcal{O}_B=\{0,\dots,2^{b_O}-1\}
\]
(with $\mathcal{O}_B=\{0\}$ when $b_O=0$), and
\[
\mathcal{R}_B:=\{r\in\mathbb{Z}:r_{\min}\le r\le r_{\max}\},
\qquad
\mathcal{R}_B^{\mathrm{enc}}:=\{0,\dots,2^{b_R}-1\}.
\]
Initialization must verify cardinality consistency
\[
|\mathcal{A}_B|=\texttt{agent\_actions}
\]
after compiling the action alphabet from $B$; on mismatch, initialization aborts
with configuration error \texttt{action\_alphabet\_mismatch}.

Let $\Theta_H$ be the complete \emph{finite} auxiliary-state set containing
exactly the remaining variables, if any, that influence future planner-facing
dynamics and are not functions of proposal source and incumbent alone.

\begin{assumption}[No hidden state outside chosen state tuple]\label{ass:no-hidden}
Every implementation variable that influences future emitted observation,
emitted reward, or decoded candidate is either:
\begin{enumerate}
    \item included in $\Theta_H$, or
    \item planner-facing inert: changing it while holding the chosen state tuple
    fixed leaves all future emitted transitions/rewards/observations unchanged.
\end{enumerate}
\end{assumption}

Define planner state space
\[
\mathcal{S}_H:=
\mathcal{Z}_H^{\mathrm{can}}\times \mathcal{Z}_{H,\mathrm{dep}}^{\mathrm{can}}\times \Theta_H,
\]
with state $s=(u,b,\theta)$ where $u$ is proposal source, $b$ incumbent,
$\theta$ auxiliary state.

Action decoding is the total deterministic map
\[
\mu_B^{\mathrm{tot}}:\mathcal{A}_B\times\mathcal{S}_H\to
\mathcal{Z}_H^{\mathrm{can}}\times\{0,1\},
\qquad
(a,s)\mapsto(\tilde z,\iota),
\]
followed by canonicalization
\[
z(a,s):=\operatorname{canon}(\tilde z).
\]
Canonical inapplicable-action semantics are invalid-without-evaluation:
when $\iota=1$, no evaluator execution is performed and no compression run is
started.

Let
\[
y_H(a,s):=
\begin{cases}
\bigl(\mathrm{invalid},\bot_c,\bot_\ell,\bot_\tau,\bot_\mu,d_{\mathrm{inapp}}\bigr),
& \iota=1,\\[3pt]
M_H^\circ\bigl(z(a,s)\bigr),
& \iota=0.
\end{cases}
\]
Define exact-objective feasibility predicate
\[
\mathrm{Feas}_H(z):=\mathbf{1}\{\mathrm{status}_H(z)=\mathrm{success}\ \land\ \theta_H(z)\ge\theta_{\min}\ \land\ \mu_H(z)\le\mu_{\max}\}.
\]
Define incumbent update
\[
\operatorname{Inc}_H(s,a):=
\begin{cases}
z(a,s),
& \iota=0\ \land\ \mathrm{Feas}_H(z(a,s))=1
\ \land\
K_H(z(a,s))\prec_K K_H(b),\\[3pt]
b,
& \text{otherwise.}
\end{cases}
\]

Let
\[
U_H:\mathcal{S}_H\times\mathcal{A}_B\to\mathcal{S}_H
\]
be the deterministic next-state map, required to satisfy
\[
\pi_b\bigl(U_H(s,a)\bigr)=\operatorname{Inc}_H(s,a),
\]
where $\pi_b(u,b,\theta)=b$ extracts the incumbent component.

\subsection{Objective-aligned reward}

Define planner-semantic reward by incumbent objective decrease:
\[
r_H(s,a):=\mathcal{J}_H\!\bigl(\pi_b(s);x\bigr)-\mathcal{J}_H\!\bigl(\pi_b(U_H(s,a));x\bigr).
\]
On any realized planner trajectory with $b_t=\pi_b(s_t)$, this agrees
stepwise with Section~\ref{sec:reward-observation}'s reward
$R_t^{\mathcal{J}}$.
This is finite and well-defined on reachable trajectories because incumbent
components are always deployable.

\begin{proposition}[Undiscounted telescoping identity]\label{prop:telescoping}
For every finite trajectory
\[
s_0,a_0,s_1,a_1,\dots,s_m,
\qquad
s_{t+1}=U_H(s_t,a_t),
\]
we have
\[
\sum_{t=0}^{m-1} r_H(s_t,a_t)
=
\mathcal{J}_H\!\bigl(\pi_b(s_0);x\bigr)-\mathcal{J}_H\!\bigl(\pi_b(s_m);x\bigr).
\]
\end{proposition}

\begin{proof}
Immediate from summing
\[
r_H(s_t,a_t)=\mathcal{J}_H\!\bigl(\pi_b(s_t);x\bigr)-\mathcal{J}_H\!\bigl(\pi_b(s_{t+1});x\bigr).
\]
\end{proof}

\subsection{Observation encoding and exact observed-Markov path}

\begin{proposition}[Baseline-derived reachable reward bound]
\label{prop:reward-bound}
Under Assumption~\ref{ass:baseline-succ}, every reachable planner-semantic
reward satisfies
\[
0\le r_H(s,a)\le \overline R_H,
\qquad
\overline R_H:=\mathcal{J}_H(b_0;x).
\]
\end{proposition}

\begin{proof}
Along every reachable trajectory from $s_0$, incumbent updates never worsen the
key and therefore never increase the incumbent objective. Hence for every
reachable state $s$,
\[
0\le \mathcal{J}_H\!\bigl(\pi_b(s);x\bigr)\le \mathcal{J}_H(b_0;x)=\overline R_H.
\]
Also,
\[
\mathcal{J}_H\!\bigl(\pi_b(U_H(s,a));x\bigr)
\le
\mathcal{J}_H\!\bigl(\pi_b(s);x\bigr),
\]
so
\[
0
\le
r_H(s,a)
=
\mathcal{J}_H\!\bigl(\pi_b(s);x\bigr)-\mathcal{J}_H\!\bigl(\pi_b(U_H(s,a));x\bigr)
\le
\overline R_H.
\]
\end{proof}

Choose deterministic state encoder
\[
\Psi_H:\mathcal{S}_H\to\mathcal{O}_B^{L_O}.
\]

\begin{definition}[Exact-state observation path]
The embedding uses the \emph{exact-state observation path} if $\Psi_H$ is
injective.
\end{definition}

A sufficient cardinality condition is
\[
2^{b_OL_O}\ge |\mathcal{S}_H|.
\]

\begin{remark}[Normative reward-encoding safety check]
\label{rem:reward-encoding-safety}
Initialization must verify
\[
r_{\min}\le 0,
\qquad
r_{\max}\ge \overline R_H,
\qquad
\rho_{\mathrm{off}}=-r_{\min},
\qquad
r_{\max}-r_{\min}\le 2^{b_R}-1.
\]
If this check fails, initialization aborts with configuration error
\texttt{reward\_encoding\_unsafe}. By
Proposition~\ref{prop:reward-bound}, this implies the finite realized reward
set satisfies
\[
\{r_H(s,a):(s,a)\in\mathcal{S}_H\times\mathcal{A}_B\}\subseteq\mathcal{R}_B,
\]
and that $Q_R$ is exact on all reachable rewards.
\end{remark}

Define exact reward encoder
\[
Q_R(r):=r+\rho_{\mathrm{off}},
\]
with admissibility condition
\[
0\le Q_R(r)\le 2^{b_R}-1
\quad\text{for all }r\in\mathcal{R}_B.
\]

The emitted percept is
\[
x_{t+1}^{\mathrm{enc}}=
\bigl(\Psi_H(s_{t+1}),Q_R(r_H(s_t,a_t))\bigr)
\in\mathcal{O}_B^{L_O}\times\mathcal{R}_B^{\mathrm{enc}}.
\]

\begin{proposition}[Finite MDP by construction]\label{prop:finite-mdp}
Under Assumptions~\ref{ass:finite-Z}--\ref{ass:no-hidden}, the induced
planner-facing process is a finite time-homogeneous MDP
\[
\mathfrak{M}_H:=(\mathcal{S}_H,\mathcal{A}_B,T_H,R_H),
\]
with deterministic transition kernel
\[
T_H(s'\mid s,a):=\mathbf{1}\{U_H(s,a)=s'\},
\]
and reward function
\[
R_H(s,a,s'):=r_H(s,a).
\]
\end{proposition}

\begin{proof}
By Assumption~\ref{ass:no-hidden}, every variable affecting future
planner-facing emissions is either part of $s$ or planner-facing inert. Hence
$s_{t+1}$ and $r_t$ are functions only of $(s_t,a_t)$. Since
$\mathcal{S}_H$ and $\mathcal{A}_B$ are finite, the process is a finite
time-homogeneous MDP.
\end{proof}

\begin{proposition}[Exact observed $1$-Markov path]\label{prop:obs-1-markov}
If the exact-state observation path is used, then the emitted percept process is
observed $1$-Markov:
\[
\Pr\!\bigl(x_{t+1}^{\mathrm{enc}}\mid a_t,x_{1:t}^{\mathrm{enc}},a_{1:t-1}\bigr)
=
\Pr\!\bigl(x_{t+1}^{\mathrm{enc}}\mid a_t,x_t^{\mathrm{enc}}\bigr).
\]
\end{proposition}

\begin{proof}
Injectivity of $\Psi_H$ implies $x_t^{\mathrm{enc}}$ determines $s_t$
uniquely. Apply Proposition~\ref{prop:finite-mdp}.
\end{proof}

\subsection{Budget-free finite-horizon planning variant}

In the budget-free variant, no countdown or terminal flag is part of
environment state. The planning horizon $m\in\mathbb{N}$ is an external planner
parameter. Proposition~\ref{prop:finite-mdp} then yields a finite
(time-homogeneous) MDP for every fixed planning horizon $m$.

\subsection{Conditional planner convergence on exact finite-horizon MC-AIXI path}\label{sec:planner-convergence-mcaixi}

Define the globally $K_H$-optimal deployable candidate
\[
z_H^{\star}:=\arg\min_{z\in\mathcal{Z}_{H,\mathrm{dep}}^{\mathrm{can}}}K_H(z),
\]
which is unique by finiteness and total ordering of $K_H$.

For fixed horizon $m\in\mathbb{N}$ and action sequence
$a_{0:m-1}\in\mathcal{A}_B^m$, define the deterministic rollout
\[
s_0^{a}:=s_0,
\qquad
s_{t+1}^{a}:=U_H(s_t^{a},a_t),
\quad t=0,\dots,m-1,
\]
and deterministic undiscounted return
\[
G_m(s_0,a_{0:m-1})
:=
\sum_{t=0}^{m-1} r_H(s_t^{a},a_t).
\]

\begin{assumption}[Planner-convergence regime]\label{ass:planner-conv-regime}
The planner is MC-AIXI(FAC-CTW) with $\rho$UCT on the exact model
$\mathfrak{M}_H=(\mathcal{S}_H,\mathcal{A}_B,T_H,R_H)$,
using objective-aligned reward $r_H$, deterministic undiscounted
finite-horizon return $G_m$,
and the budget-free finite-horizon variant of this section.
\end{assumption}

\begin{assumption}[Horizon-$m$ attainability with unique optimal action path]\label{ass:planner-conv-unique-path}
There exists $m\in\mathbb{N}$ and a unique action sequence
$a_{0:m-1}^{\star}\in\mathcal{A}_B^m$. Define optimal-prefix states by
\[
s_0^{\star}:=s_0,
\qquad
s_{t+1}^{\star}:=U_H(s_t^{\star},a_t^{\star}),
\quad t=0,\dots,m-1.
\]
Then:
\begin{enumerate}
    \item $\pi_b(s_m^{\star})=z_H^{\star}$, and
    \item for every other sequence $a_{0:m-1}\neq a_{0:m-1}^{\star}$,
    \[
    G_m(s_0,a_{0:m-1}^{\star}) > G_m(s_0,a_{0:m-1}),
    \]
\end{enumerate}
\end{assumption}

\begin{assumption}[Per-decision asymptotic root-action consistency]\label{ass:planner-conv-root-consistency}
For each simulation budget $N\in\mathbb{N}$, let
$S_t^{(N)}\in\mathcal{S}_H$ be the closed-loop state process defined by
\[
S_0^{(N)}:=s_0,
\qquad
S_{t+1}^{(N)}:=U_H\big(S_t^{(N)},a_t^{(N)}\big),
\]
where $a_t^{(N)}$ is the root action selected at step $t$ using $N$
planner-internal simulations. For each $t\in\{0,\dots,m-1\}$,
\[
\mathbb{P}\!\big(a_t^{(N)}=a_t^{\star}\mid S_t^{(N)}=s_t^{\star}\big)
\xrightarrow[N\to\infty]{}1.
\]
\end{assumption}

\begin{theorem}[Conditional MC-AIXI convergence to $z_H^{\star}$]\label{thm:planner-convergence-mcaixi}
Under Assumptions~\ref{ass:planner-conv-regime}--\ref{ass:planner-conv-root-consistency}
and Proposition~\ref{prop:finite-mdp},
consider the closed-loop process indexed by simulation budget $N$, and define
$b_t^{(N)}:=\pi_b\big(S_t^{(N)}\big)$. Then
\[
\mathbb{P}\!\big(b_m^{(N)}=z_H^{\star}\big)\xrightarrow[N\to\infty]{}1.
\]
\end{theorem}

For the non-unique maximizer case, see Proposition~\ref{prop:planner-convergence-optset}.

\begin{proof}
We prove by induction that
\[
\mathbb{P}\!\big(S_t^{(N)}=s_t^{\star}\big)\xrightarrow[N\to\infty]{}1
\quad\text{for each }t=0,\dots,m.
\]
Base case $t=0$ is immediate from $S_0^{(N)}=s_0=s_0^{\star}$.

Induction step: assume
$\mathbb{P}(S_t^{(N)}=s_t^{\star})\to 1$ for some $t<m$.
Then
\[
\begin{aligned}
\mathbb{P}\!\big(S_{t+1}^{(N)}=s_{t+1}^{\star}\big)
&\ge
\mathbb{P}\!\big(S_t^{(N)}=s_t^{\star},\ a_t^{(N)}=a_t^{\star}\big)\\
&=
\mathbb{P}\!\big(S_t^{(N)}=s_t^{\star}\big)
\,\mathbb{P}\!\big(a_t^{(N)}=a_t^{\star}\mid S_t^{(N)}=s_t^{\star}\big).
\end{aligned}
\]
The first factor tends to $1$ by induction hypothesis and the second by
Assumption~\ref{ass:planner-conv-root-consistency}, so
\[
\liminf_{N\to\infty}\mathbb{P}\!\big(S_{t+1}^{(N)}=s_{t+1}^{\star}\big)\ge 1.
\]
Hence
$\mathbb{P}(S_{t+1}^{(N)}=s_{t+1}^{\star})\to 1$.

By induction,
$\mathbb{P}(S_m^{(N)}=s_m^{\star})\to 1$.
From Assumption~\ref{ass:planner-conv-unique-path}, $\pi_b(s_m^{\star})=z_H^{\star}$, so
\[
\mathbb{P}\!\big(b_m^{(N)}=z_H^{\star}\big)
\ge
\mathbb{P}\!\big(S_m^{(N)}=s_m^{\star}\big)
\xrightarrow[N\to\infty]{}1,
\]
which proves the claim.
\end{proof}

\begin{proposition}[Optimal-action-set variant (non-unique optimal sequences)]\label{prop:planner-convergence-optset}
In the setting of Theorem~\ref{thm:planner-convergence-mcaixi}, replace
Assumption~\ref{ass:planner-conv-unique-path} by the following:
\begin{enumerate}
    \item define
    \[
    \mathcal{A}_m^{\star}:=
    \arg\max_{a_{0:m-1}\in\mathcal{A}_B^m}G_m(s_0,a_{0:m-1}),
    \]
    which may be non-singleton;
    \item for each sequence $a_{0:m-1}$ define deterministic trajectory
    \[
    s_0^{a}:=s_0,
    \qquad
    s_{t+1}^{a}:=U_H(s_t^{a},a_t),
    \quad t=0,\dots,m-1,
    \]
    and require
    \[
    \pi_b(s_m^{a})=z_H^{\star}
    \quad\text{for all }a_{0:m-1}\in\mathcal{A}_m^{\star}.
    \]
\end{enumerate}
For each $t\in\{0,\dots,m\}$ define
\[
\mathcal{S}_t^{\star}:=\{s_t^{a}:a_{0:m-1}\in\mathcal{A}_m^{\star}\}.
\]
For each $t\in\{0,\dots,m-1\}$ and $s\in\mathcal{S}_t^{\star}$ define
\[
\mathcal{U}_t^{\star}(s):=\{a\in\mathcal{A}_B:U_H(s,a)\in\mathcal{S}_{t+1}^{\star}\}.
\]
Assume also that for each $t\in\{0,\dots,m-1\}$,
\[
\inf_{s\in\mathcal{S}_t^{\star}}
\mathbb{P}\!\big(a_t^{(N)}\in\mathcal{U}_t^{\star}(s)\mid S_t^{(N)}=s\big)
\xrightarrow[N\to\infty]{}1.
\]
Then
\[
\mathbb{P}\!\big(b_m^{(N)}=z_H^{\star}\big)\xrightarrow[N\to\infty]{}1.
\]
\end{proposition}

\begin{proof}
Let
\[
F_t^{(N)}:=\{S_t^{(N)}\in\mathcal{S}_t^{\star}\},
\qquad t=0,\dots,m.
\]
Since $S_0^{(N)}=s_0$ and $s_0\in\mathcal{S}_0^{\star}$,
$\mathbb{P}(F_0^{(N)})=1$.

If $F_t^{(N)}$ occurs and
$a_t^{(N)}\in\mathcal{U}_t^{\star}(S_t^{(N)})$, then by definition of
$\mathcal{U}_t^{\star}$ we have $S_{t+1}^{(N)}\in\mathcal{S}_{t+1}^{\star}$,
hence
\[
\mathbb{P}(F_{t+1}^{(N)})
\ge
\mathbb{P}(F_t^{(N)})
\inf_{s\in\mathcal{S}_t^{\star}}
\mathbb{P}\!\big(a_t^{(N)}\in\mathcal{U}_t^{\star}(s)\mid S_t^{(N)}=s\big).
\]
By the assumed uniform conditional consistency, the multiplicative factor tends
to $1$ for each fixed $t$, so induction on $t$ gives
$\mathbb{P}(F_m^{(N)})\to 1$.

By common-terminal-incumbent optimality,
$\pi_b(s)=z_H^{\star}$ for every $s\in\mathcal{S}_m^{\star}$, therefore
\[
\mathbb{P}\!\big(b_m^{(N)}=z_H^{\star}\big)
\ge
\mathbb{P}(F_m^{(N)})
\xrightarrow[N\to\infty]{}1.
\]
\end{proof}

\begin{remark}[Informal takeaway for planner convergence]
Assumption~\ref{ass:planner-conv-unique-path} is a stronger-than-ideal
single-path sufficient condition used for the clean theorem statement. The
broader interpretation is set-valued: what matters is asymptotic selection of
actions that remain inside the optimal action set leading to the common
terminal incumbent $z_H^{\star}$, as formalized in
Proposition~\ref{prop:planner-convergence-optset}.
\end{remark}

\begin{remark}
Theorem~\ref{thm:planner-convergence-mcaixi} and
Proposition~\ref{prop:planner-convergence-optset} are asymptotic-in-simulations
statements for the exact objective-aligned undiscounted finite-horizon
MC-AIXI path. They are not finite-budget runtime guarantees and are not, by
themselves, convergence theorems for runtime \texttt{aiqi\_discounted} or
runtime \texttt{aiqi\_warmstart\_exact\allowbreak\_jh} under finite simulation budgets
and optional exploratory/non-greedy control.
Their theorem target is $b_m^{(N)}$ (incumbent after exactly $m$ executed planner
steps), not the runtime output
$z^\star=\arg\min_{z\in\mathcal{S}_{\mathrm{term}}^{\mathrm{dep}}}K_H(z)$ from
Section~\ref{sec:aixi-runtime} under wall-clock-bounded termination.
\end{remark}

\subsection{Optional budgeted/episodic variant}

To model explicit per-episode bounds, add deterministic countdown component
\[
d_t\in\{0,\dots,D_{\max}\}
\]
to $\Theta_H$ and update it inside $U_H$. The Markov property is preserved because the
countdown is inside state.

\subsection{Optional reset-augmented continuing variant}

Let $\mathsf{Term}:\mathcal{S}_H\to\{0,1\}$ be a terminal predicate and $\nu$ a
restart distribution on $\mathcal{S}_H$. Define
\[
\widetilde T_H(s'\mid s,a):=
\begin{cases}
T_H(s'\mid s,a), & \mathsf{Term}(s)=0,\\[3pt]
\nu(s'), & \mathsf{Term}(s)=1.
\end{cases}
\]
Optionally define lazy kernel
\[
\widetilde T_H^{(\lambda)}:=\lambda I+(1-\lambda)\widetilde T_H,
\qquad
\lambda\in(0,1).
\]
For any fixed action-selection policy, if the induced finite chain is
irreducible then the lazy chain is aperiodic and therefore ergodic.

\subsection{Runtime controller path for MC-AIXI and AIQI-family controllers}\label{sec:aixi-runtime}

This subsection is implementation-normative: it specifies runtime execution
when \texttt{algorithm} is \texttt{mc\_aixi\_fac\_ctw},
\texttt{aiqi\_discounted}, or
\texttt{aiqi\_warmstart\_exact\allowbreak\_jh}.

Given the same input contract and evaluator definitions used by
Section~\ref{sec:ahc-runtime}, define
\[
N_{\mathrm{sim}}^{\max}:=\texttt{planner\_simulations\_per\_step}\in\mathbb{N},
\qquad N_{\mathrm{sim}}^{\max}\ge 1,
\]
and apply the following execution contract.

\paragraph{Shared initialization contract.}
\begin{enumerate}
    \item compile planner action alphabet $\mathcal{A}_B$ from $B$ and verify
    cardinality consistency
    \[
    |\mathcal{A}_B|=\texttt{agent\_actions};
    \]
    on mismatch, abort initialization with configuration error
    \texttt{action\_alphabet\_mismatch},
    \item after the normative baseline evaluation, verify the exact
    reward-encoding safety check from
    Remark~\ref{rem:reward-encoding-safety}; on failure, abort initialization
    with configuration error \texttt{reward\_encoding\allowbreak\_unsafe},
    \item instantiate planner-facing environment state $s_0\in\mathcal{S}_H$
    with the post-warmup baseline incumbent satisfying the normative baseline
    precondition from Section~\ref{subsec:baseline-precondition}
    (equivalently Assumption~\ref{ass:baseline-succ} in the fixed-profile
    analysis).
\end{enumerate}

\paragraph{Controller-specific internal policy contract.}
Initialize controller-specific internal state by
\[
\begin{aligned}
u_0^{\mathrm{mc}}&\ \text{(MC-AIXI/FAC-CTW)},\\
u_0^{\mathrm{aiqi,disc}}&\ \text{(discounted variant)},\\
u_0^{\mathrm{aiqi,warm}}&\ \text{(warm-start exact variant)}.
\end{aligned}
\]
Apply seeded randomized tie-breaking for MC-AIXI; apply deterministic
lowest-index greedy tie-breaking for \texttt{aiqi\_discounted}
(Section~\ref{sec:aiqi-discounted-contract}) and for the exact-greedy
warm-start path (Section~\ref{sec:aiqi-warmstart-exact-jh}); explicit
exploration parameters are heuristic. At each decision step $t$, allocate
planner-internal simulation count
\[
N_{\mathrm{sim},t}\in\{1,\dots,N_{\mathrm{sim}}^{\max}\}
\]
as permitted by remaining wall-clock budget, run those internal simulations,
and then select external action $a_t\in\mathcal{A}_B$ from the current
controller state.

For \texttt{algorithm = aiqi\_warmstart\_exact\allowbreak\_jh} with
$R_{\mathrm{si}}>1$, this single-pass execution is used as the per-round
kernel inside the bounded same-task trace-refresh wrapper of
Section~\ref{sec:aiqi-warmstart-self-improve}. The global wall-clock budget
remains $T_{\mathrm{total}}$, the next round starts from the current outer
incumbent, and the final output is still the minimum-$K_H$ deployable candidate
over all real evaluations performed across all rounds.

\paragraph{Shared environment, termination, and output contract.}
\begin{enumerate}
    \item environment applies deterministic transition
    $s_{t+1}=U_H(s_t,a_t)$ and emits encoded percept
    \[
    x_{t+1}^{\mathrm{enc}}=
    \bigl(\Psi_H(s_{t+1}),Q_R(r_H(s_t,a_t))\bigr),
    \]
    with canonical inapplicable-action semantics inherited from
    Section~\ref{sec:aixi-embedding}: when decoder flag $\iota=1$, the emitted
    outcome is invalid with diagnostic token
    \texttt{inapplicable\_action}, no evaluator execution occurs, and no
    compression run is started,
    \item planner updates its internal state using
    $(a_t,x_{t+1}^{\mathrm{enc}})$ and repeats until termination criteria are
    met,
    \item termination criteria are the same normative criteria used by the
    tuning runtime (global deadline, \texttt{max\_evaluations} if set,
    neighborhood exhaustion when provable, unrecoverable evaluator failure),
    \item final output candidate is
    \[
    z^\star=\arg\min_{z\in\mathcal{S}_{\mathrm{term}}^{\mathrm{dep}}}K_H(z),
    \]
    where
    \[
    \mathcal{S}_{\mathrm{term}}^{\mathrm{dep}}:=
    \left\{z\in\mathcal{Z}_H^{\mathrm{can}}:\
    \begin{aligned}[t]
    &z\text{ evaluated before termination},\\
    &\mathrm{status}_H(z)=\mathrm{success},\\
    &\theta_H(z)\ge\theta_{\min},\ \mu_H(z)\le\mu_{\max}
    \end{aligned}
    \right\}.
    \]
\end{enumerate}
Controller-specific search-control dynamics may differ, but scoring/output
semantics and candidate-validity semantics remain shared.

\subsection{Normative discounted AIQI controller contract}
\label{sec:aiqi-discounted-contract}

This subsection is normative for
\texttt{algorithm = aiqi\_discounted}. It defines a
discounted/bin-quantized heuristic controller path and does \emph{not} provide
exact-$\mathcal{J}_H$ guarantees for finite-horizon objective alignment.

Define controller parameters
\[
H_{\mathrm{ret}}:=\texttt{return\_horizon}\in\mathbb{N},\ H_{\mathrm{ret}}\ge 1,
\qquad
\gamma:=\texttt{discount\_factor}\in[0,1),
\]
\[
M:=\texttt{return\_bins}\in\mathbb{N},
\qquad
M\text{ is a power of two}.
\]

Define clamp on unit interval
\[
\operatorname{clamp}_{[0,1]}(x):=\min\!\bigl(\max(x,0),1\bigr).
\]
For semantic rewards $r_t\in\mathcal{R}_B$, define normalized clipped reward
\[
\hat r_t:=\operatorname{clamp}_{[0,1]}\!\left(
\frac{r_t-r_{\min}}{r_{\max}-r_{\min}}
\right),
\]
and discounted finite-horizon target
\[
G_{t,H_{\mathrm{ret}}}^{(\gamma)}:=
\operatorname{clamp}_{[0,1]}\!\left(
(1-\gamma)\sum_{k=0}^{H_{\mathrm{ret}}-1}\gamma^k\hat r_{t+k}
\right).
\]

Define return-bin quantizer
\[
Q_{\mathrm{return}}(G):=
\min\!\left(\left\lfloor M G\right\rfloor,\ M-1\right),
\qquad
G\in[0,1].
\]
The discounted controller predicts/optimizes the distribution of
$Q_{\mathrm{return}}(G_{t,H_{\mathrm{ret}}}^{(\gamma)})$, not exact
$\mathcal{J}_H$-improvement targets.

Greedy action ties are broken deterministically by lowest action index; any
explicit exploration parameter is heuristic and outside exact-$\mathcal{J}_H$
claims.

\subsection{MC-AIXI-warm-started model-free exact-\texorpdfstring{$\mathcal{J}_H$}{J\_H} controller}
\label{sec:aiqi-warmstart-exact-jh}

This subsection defines a model-free controller path whose $m$-step prediction
target is exactly aligned with the deployable tuning objective $\mathcal{J}_H$
via terminal-incumbent improvement, while using MC-AIXI-generated teacher data
to initialize the label predictor. It is inspired by the delayed-label
phase-predictor structure of AIQI, but it does \emph{not} use discounted return
targets and does \emph{not} invoke the discounted AIQI theorem directly.
It is normative for \texttt{algorithm = aiqi\_warmstart\_exact\allowbreak\_jh}.

\paragraph{Setup.}
Fix the same planner-facing finite interaction contract as in
Section~\ref{sec:reward-observation}: finite action alphabet
$\mathcal{A}_B$, finite encoded observation alphabet, deterministic incumbent
key, and the objective-aligned reward
\[
R_t^{\mathcal{J}}:=\mathcal{J}_H(b_t;x)-\mathcal{J}_H(b_{t+1};x).
\]
Define
\[
m:=\texttt{return\_horizon}\in\mathbb{N},\qquad m\ge 1.
\]

Define the exact $m$-step cumulative improvement target
\[
G_t^{(m)}
:=
\sum_{k=0}^{m-1} R_{t+k}^{\mathcal{J}}.
\]
Because the reward is incumbent-objective difference, this telescopes exactly:
\[
G_t^{(m)}
=
\mathcal{J}_H(b_t;x)-\mathcal{J}_H(b_{t+m};x).
\]
Hence maximizing expected $G_t^{(m)}$ is equivalent to minimizing the expected
terminal incumbent objective after $m$ further planner steps.

\begin{remark}[Integer-valued exact target]
Operationally, $\mathcal{J}_H$ is byte-valued, so
$R_t^{\mathcal{J}}\in\mathbb{Z}_{\ge 0}$ and
$G_t^{(m)}\in\mathbb{Z}_{\ge 0}$.
A safe finite label alphabet is therefore
\[
\mathcal{G}_m:=\{0,1,\dots,\overline G_m\},
\qquad
\overline G_m:=\mathcal{J}_H(b_0;x),
\]
since $\mathcal{J}_H(b_{t+m};x)\ge 0$ and incumbents never worsen.
Any tighter finite upper bound may be used instead.
\end{remark}

\paragraph{Initial teacher dataset and admissible same-task trace sources.}
Let
\[
\mathcal{D}_{\mathrm{teach}}
=
\bigl\{
(h_{<t}^{(i)},a_t^{(i)},G_t^{(m),(i)})
\bigr\}_{i=1}^{N_{\mathrm{teach}}}
\]
be the initial teacher dataset.

\begin{definition}[Admissible same-task teacher dataset]
\label{def:admissible-teacher}
Fix the exact-$\mathcal{J}_H$ interaction contract of this subsection.
A finite dataset $\mathcal{D}_{\mathrm{teach}}$ is \emph{admissible} if every
record is collected from a realized controller trajectory on the same
$(B,H,x)$-induced tuning environment under the same finite planner-visible
history/action contract, and each label is computed from the realized
incumbent path by
\[
G_t^{(m),(i)}
=
\mathcal{J}_H(b_t^{(i)};x)-\mathcal{J}_H(b_{t+m}^{(i)};x).
\]
\end{definition}

MC-AIXI executed trajectories are the canonical source of an initial admissible
dataset. Prior executions of the exact-$\mathcal{J}_H$ controller are also
admissible. More generally, any controller can supply teacher traces exactly
when its realized trajectory is exposed through the same action/history
contract and the same exact label semantics above. Operationally, the initial
admissible dataset is provided by
\texttt{warmstart\_teacher\_dataset\_path}.

\begin{remark}[Which other controllers qualify]
Annealed hill climbing or another optimizer is not directly reusable here
unless an explicit semantics-preserving adapter exposes its realized trace
through the same compiled action alphabet, same history convention, and same
exact label definition $G_t^{(m)}$. Without that contract match, its traces are
non-admissible for this controller path.
\end{remark}

\paragraph{Phase-indexed delayed-label construction.}
Define augmentation period
\[
N:=\texttt{label\_phase\_period}\in\mathbb{N},\qquad N\ge m.
\]
For each phase $n\in\{0,\dots,N-1\}$, define the phase-restricted teacher set
\[
\mathcal{D}_{\mathrm{teach}}^{(n)}
:=
\left\{
(h_{<t}^{(i)},a_t^{(i)},G_t^{(m),(i)}):
\ t \equiv n \pmod N
\right\}.
\]
For each phase $n$, define a predictive distribution
\[
\psi_n(\,\cdot\,\mid h_{<t}a_t)
\in
\Delta(\mathcal{G}_m),
\]
where $\Delta(\mathcal{G}_m)$ denotes the probability simplex on
$\mathcal{G}_m$.

The initialization (warm-start) of $\psi_n$ from
$\mathcal{D}_{\mathrm{teach}}^{(n)}$ is controller-internal and may be implemented
by any deterministic finite-alphabet predictive mechanism, provided the resulting
predictor is a well-defined conditional distribution on $\mathcal{G}_m$.
Examples include exact count tables, CTW-style pseudocount initialization on an
augmented label sequence, or any other deterministic finite-alphabet Bayesian or
MDL-compatible predictor.

Define the unified phase predictor
\[
\psi(\,\cdot\,\mid h_{<t}a_t):=
\psi_{t\bmod N}(\,\cdot\,\mid h_{<t}a_t).
\]

\paragraph{Warm-started action-value estimate.}
Define the exact-$\mathcal{J}_H$ model-free action-value estimate
\[
\widehat Q_m(h_{<t},a)
:=
\mathbb{E}_{\psi}\!\left[G_t^{(m)}\mid h_{<t},a\right]
=
\sum_{g\in\mathcal{G}_m} g\,\psi(g\mid h_{<t}a).
\]

\paragraph{True conditional law and exact action-value.}
For each history-action pair, define the true conditional law of the exact
$m$-step label by
\[
p_t(\,\cdot\,\mid h_{<t},a)
:=
\mathbb{P}\!\left(G_t^{(m)}\in\,\cdot\,\mid h_{<t},a\right)
\in
\Delta(\mathcal{G}_m),
\]
and the corresponding exact action-value by
\[
Q_t^\star(h_{<t},a)
:=
\mathbb{E}_{p_t}\!\left[G_t^{(m)}\mid h_{<t},a\right]
=
\sum_{g\in\mathcal{G}_m} g\,p_t(g\mid h_{<t},a).
\]
For $p,q\in\Delta(\mathcal{G}_m)$, write total variation distance as
\[
\|p-q\|_{\mathrm{TV}}
:=
\sup_{A\subseteq\mathcal{G}_m}|p(A)-q(A)|
=
\frac{1}{2}\sum_{g\in\mathcal{G}_m}|p(g)-q(g)|.
\]

\paragraph{Deployment policy.}
The greedy deployment rule for this controller is
\[
\hat a_t
\in
\arg\max_{a\in\mathcal{A}_B}\widehat Q_m(h_{<t},a),
\]
with deterministic tie-breaking (e.g.\ lowest-index maximizing action).

Optionally, one may use an exploration schedule
\[
\varepsilon_t\in[0,1],
\]
and select actions by
\[
\pi_t(a\mid h_{<t})
=
(1-\varepsilon_t)\mathbf{1}[a=\hat a_t]
+
\varepsilon_t\,u_{\mathcal{A}_B}(a),
\]
where $u_{\mathcal{A}_B}$ is the uniform distribution on $\mathcal{A}_B$.
However, the exact action-selection claim below applies only to the greedy case
$\varepsilon_t=0$.

\paragraph{Online update.}
After execution reveals the realized label $G_t^{(m)}$, the active phase
predictor $\psi_{t\bmod N}$ may be updated online using that label.
Online updating is optional for semantic exactness but useful for correcting
teacher-induced distribution shift.

\begin{proposition}[Exact finite-horizon objective alignment]
\label{prop:aiqi-warmstart-exact-jh}
For every history $h_{<t}$ and action $a\in\mathcal{A}_B$,
\[
\mathbb{E}\!\left[G_t^{(m)}\mid h_{<t},a\right]
=
\mathcal{J}_H(b_t;x)
-
\mathbb{E}\!\left[\mathcal{J}_H(b_{t+m};x)\mid h_{<t},a\right].
\]
Consequently,
\[
\arg\max_{a\in\mathcal{A}_B}
\mathbb{E}\!\left[G_t^{(m)}\mid h_{<t},a\right]
=
\arg\min_{a\in\mathcal{A}_B}
\mathbb{E}\!\left[\mathcal{J}_H(b_{t+m};x)\mid h_{<t},a\right].
\]
Therefore, if $\psi(\cdot\mid h_{<t}a)$ equals the true conditional law of
$G_t^{(m)}$ given $(h_{<t},a)$, then the greedy rule
$\hat a_t\in\arg\max_a \widehat Q_m(h_{<t},a)$ is exactly optimal for the
current decision with respect to the finite-horizon terminal-incumbent
objective at horizon $m$.
\end{proposition}

\begin{proof}
By definition,
\[
G_t^{(m)}
=
\sum_{k=0}^{m-1}R_{t+k}^{\mathcal{J}}
=
\sum_{k=0}^{m-1}
\left(
\mathcal{J}_H(b_{t+k};x)-\mathcal{J}_H(b_{t+k+1};x)
\right),
\]
which telescopes to
\[
G_t^{(m)}
=
\mathcal{J}_H(b_t;x)-\mathcal{J}_H(b_{t+m};x).
\]
Condition on $(h_{<t},a)$ and take expectations. Since
$\mathcal{J}_H(b_t;x)$ is fixed given $h_{<t}$, maximizing the expected
cumulative improvement is equivalent to minimizing the expected terminal
objective at horizon $m$.
\end{proof}

\begin{theorem}[Approximate finite-horizon greedy optimality under total-variation predictor error]
\label{thm:aiqi-warmstart-tv-robust}
Fix a history $h_{<t}$ and suppose
\[
\sup_{a\in\mathcal{A}_B}
\left\|
\psi(\cdot\mid h_{<t}a)-p_t(\cdot\mid h_{<t},a)
\right\|_{\mathrm{TV}}
\le \delta_t.
\]
Let
\[
a_t^\star\in\arg\max_{a\in\mathcal{A}_B}Q_t^\star(h_{<t},a),
\qquad
\hat a_t\in\arg\max_{a\in\mathcal{A}_B}\widehat Q_m(h_{<t},a),
\]
where $\hat a_t$ uses the deterministic tie-breaking rule above. Then
\[
\sup_{a\in\mathcal{A}_B}
\left|
\widehat Q_m(h_{<t},a)-Q_t^\star(h_{<t},a)
\right|
\le
\overline G_m\,\delta_t,
\]
and
\[
Q_t^\star(h_{<t},a_t^\star)-Q_t^\star(h_{<t},\hat a_t)
\le
2\overline G_m\,\delta_t.
\]
Equivalently,
\[
\mathbb{E}\!\left[\mathcal{J}_H(b_{t+m};x)\mid h_{<t},\hat a_t\right]
-
\min_{a\in\mathcal{A}_B}
\mathbb{E}\!\left[\mathcal{J}_H(b_{t+m};x)\mid h_{<t},a\right]
\le
2\overline G_m\,\delta_t.
\]
\end{theorem}

\begin{proof}
Fix any action $a\in\mathcal{A}_B$. Since
$G_t^{(m)}\in\mathcal{G}_m\subseteq[0,\overline G_m]$, the bounded-function
characterization of total variation on a finite alphabet gives
\[
\begin{aligned}
\left|
\widehat Q_m(h_{<t},a)-Q_t^\star(h_{<t},a)
\right|
&=
\left|
\mathbb{E}_{\psi}\!\left[G_t^{(m)}\mid h_{<t},a\right]
-
\mathbb{E}_{p_t}\!\left[G_t^{(m)}\mid h_{<t},a\right]
\right|\\
&\le
\overline G_m
\left\|
\psi(\cdot\mid h_{<t}a)-p_t(\cdot\mid h_{<t},a)
\right\|_{\mathrm{TV}}.
\end{aligned}
\]
Taking the supremum over actions yields
\[
\sup_{a\in\mathcal{A}_B}
\left|
\widehat Q_m(h_{<t},a)-Q_t^\star(h_{<t},a)
\right|
\le
\overline G_m\,\delta_t.
\]
For the greedy-action suboptimality bound, write
\[
\begin{aligned}
Q_t^\star(h_{<t},a_t^\star)-Q_t^\star(h_{<t},\hat a_t)
&=
\bigl(Q_t^\star(h_{<t},a_t^\star)-\widehat Q_m(h_{<t},a_t^\star)\bigr)\\
&\quad+
\bigl(\widehat Q_m(h_{<t},a_t^\star)-\widehat Q_m(h_{<t},\hat a_t)\bigr)\\
&\quad+
\bigl(\widehat Q_m(h_{<t},\hat a_t)-Q_t^\star(h_{<t},\hat a_t)\bigr).
\end{aligned}
\]
The middle term is nonpositive because $\hat a_t$ maximizes
$\widehat Q_m(h_{<t},\cdot)$, so
\[
Q_t^\star(h_{<t},a_t^\star)-Q_t^\star(h_{<t},\hat a_t)
\le
2\sup_{a\in\mathcal{A}_B}
\left|
\widehat Q_m(h_{<t},a)-Q_t^\star(h_{<t},a)
\right|
\le
2\overline G_m\,\delta_t.
\]
By Proposition~\ref{prop:aiqi-warmstart-exact-jh}, for every action $a$,
\[
Q_t^\star(h_{<t},a)
=
\mathcal{J}_H(b_t;x)
-
\mathbb{E}\!\left[\mathcal{J}_H(b_{t+m};x)\mid h_{<t},a\right].
\]
Substituting this identity into the previous inequality cancels the common term
$\mathcal{J}_H(b_t;x)$ and yields the terminal-objective bound.
\end{proof}

\begin{corollary}[Exact action identification under a true action gap]
\label{cor:aiqi-warmstart-gap}
Assume $|\mathcal{A}_B|\ge 2$, and let $a_t^\star$ be the deterministic
lowest-index maximizer of $Q_t^\star(h_{<t},\cdot)$. Define the true action gap by
\[
\Delta_t(h_{<t})
:=
Q_t^\star(h_{<t},a_t^\star)
-
\max_{a\in\mathcal{A}_B\setminus\{a_t^\star\}}Q_t^\star(h_{<t},a).
\]
If
\[
\sup_{a\in\mathcal{A}_B}
\left|
\widehat Q_m(h_{<t},a)-Q_t^\star(h_{<t},a)
\right|
<
\frac{\Delta_t(h_{<t})}{2},
\]
then $\hat a_t=a_t^\star$.
\end{corollary}

\begin{proof}
For any action $a\neq a_t^\star$,
\[
\begin{aligned}
\widehat Q_m(h_{<t},a_t^\star)-\widehat Q_m(h_{<t},a)
&\ge
Q_t^\star(h_{<t},a_t^\star)-Q_t^\star(h_{<t},a)\\
&\quad-
2\sup_{a'\in\mathcal{A}_B}
\left|
\widehat Q_m(h_{<t},a')-Q_t^\star(h_{<t},a')
\right|\\
&>
\Delta_t(h_{<t})-2\cdot\frac{\Delta_t(h_{<t})}{2}
=0.
\end{aligned}
\]
Hence $\widehat Q_m(h_{<t},a_t^\star)>\widehat Q_m(h_{<t},a)$ for every
$a\neq a_t^\star$, so $a_t^\star$ is the unique maximizer of
$\widehat Q_m(h_{<t},\cdot)$. Therefore the deterministic greedy rule selects
$\hat a_t=a_t^\star$.
\end{proof}

\begin{remark}[Scope of the approximate-law result]
Theorem~\ref{thm:aiqi-warmstart-tv-robust} and
Corollary~\ref{cor:aiqi-warmstart-gap} are local finite-horizon robustness
statements. They do not constitute a long-run convergence theorem for the
runtime warm-start exact-$\mathcal{J}_H$ controller under finite simulations or
exploration.
\end{remark}

\begin{remark}[Role of admissible teacher data]
Initial MC-AIXI traces, prior exact-$\mathcal{J}_H$ reruns, or other
admissible same-task traces do not change the objective. Their role is only to
initialize or refine $\psi$ away from the cold-start regime. Thus the
controller remains model-free at deployment time, even when some teacher traces
were generated by a model-based controller.
\end{remark}

\begin{remark}[Claim boundary]
This subsection gives an exact finite-horizon objective-alignment statement for
the current decision step, not a discounted-return asymptotic
$\varepsilon$-optimality statement. If $\varepsilon_t>0$, then the executed
policy is exploratory and therefore is not exactly greedy at that decision
step. If one wants the executed action to coincide with the greedy rule
analyzed here, use $\varepsilon_t=0$ at deployment; exact optimality at that
step additionally requires the predictor-equality premise in
Proposition~\ref{prop:aiqi-warmstart-exact-jh}.
\end{remark}

\begin{remark}[Implementation-conformance split]
Conforming implementations must keep the controller kinds
\texttt{aiqi\_discounted} and \texttt{aiqi\_warmstart\_exact\allowbreak\_jh}
distinct. Reusing discounted return-bin targets under the exact warm-start
controller key is non-conforming to this specification. Likewise, any bounded
self-improvement trace-refresh wrapper may merge only admissible same-task
exact-label traces as defined above; reusing differently encoded or surrogate-
labeled traces is non-conforming.
\end{remark}

\subsection{Optional bounded same-task trace-refresh self-improvement loop}
\label{sec:aiqi-warmstart-self-improve}

This subsection defines an optional bounded wrapper around
\texttt{aiqi\_warmstart\_exact\allowbreak\_jh}. It reuses the same exact
label mechanism on the original tuning task; it does \emph{not} introduce any
auxiliary objective.

Let
\[
R_{\mathrm{si}}:=\texttt{self\_improvement\_rounds}\in\mathbb{N},
\qquad
R_{\mathrm{si}}\ge 1,
\]
with default $R_{\mathrm{si}}=1$.
When $R_{\mathrm{si}}>1$, interpret $T_{\mathrm{total}}$ as the single
outer-loop wall-clock budget and define deterministic round deadlines
\[
\Delta_r^{\mathrm{si}}:=\frac{r}{R_{\mathrm{si}}}T_{\mathrm{total}},
\qquad r=0,\dots,R_{\mathrm{si}}.
\]

Let the round-$r$ teacher dataset be $\mathcal{D}^{(r)}$, with initial dataset
$\mathcal{D}^{(0)}:=\mathcal{D}_{\mathrm{teach}}$, and define
\[
\beta_0:=b_0,
\qquad
\mathcal{S}_0^{\mathrm{meta}}:=\{\beta_0\}.
\]
For each round $r\in\{0,\dots,R_{\mathrm{si}}-1\}$:
\begin{enumerate}
    \item initialize the phase predictors from $\mathcal{D}^{(r)}$,
    \item initialize the round-$r$ controller with incumbent $\beta_r$ and
    execute the exact-$\mathcal{J}_H$ controller on the same $(B,H,x)$ task
    until the earlier of global termination and wall-clock deadline
    $\Delta_{r+1}^{\mathrm{si}}$,
    \item let $L_r\in\mathbb{N}_0$ be the number of executed actions in round
    $r$, and define the realized labeled round-trace set
    \[
    \mathcal{T}^{(r)}
    :=
    \left\{
    (h_{<t}^{(r)},a_t^{(r)},G_t^{(m),(r)}):\ 0\le t\le L_r-m
    \right\},
    \]
    with $\mathcal{T}^{(r)}:=\varnothing$ when $L_r<m$,
    \item let $\mathcal{E}_r^{\mathrm{dep}}$ be the set of deployable
    candidates actually evaluated in round $r$, and update
    \[
    \mathcal{D}^{(r+1)}:=\operatorname{Merge}\!\bigl(\mathcal{D}^{(r)},\mathcal{T}^{(r)}\bigr),
    \qquad
    \mathcal{S}_{r+1}^{\mathrm{meta}}:=\mathcal{S}_r^{\mathrm{meta}}\cup \mathcal{E}_r^{\mathrm{dep}},
    \]
    where $\operatorname{Merge}$ is any deterministic deduplicating union rule,
    \item define the next outer incumbent by
    \[
    \beta_{r+1}:=\arg\min_{z\in\mathcal{S}_{r+1}^{\mathrm{meta}}}K_H(z).
    \]
\end{enumerate}
At termination the wrapper returns $\beta_{R_{\mathrm{si}}}$.

\begin{proposition}[Same-task safety of bounded trace refresh]
\label{prop:aiqi-warmstart-self-improve-safe}
For the bounded same-task trace-refresh loop above:
\begin{enumerate}
    \item each round uses the same exact label $G_t^{(m)}$ and the same
    finite-horizon terminal-incumbent objective as
    Proposition~\ref{prop:aiqi-warmstart-exact-jh},
    \item every teacher record added through $\mathcal{T}^{(r)}$ is extracted
    from a realized trajectory on the original task and therefore does not
    require any auxiliary evaluation objective,
    \item the outer incumbent sequence is monotone:
    \[
    K_H(\beta_{r+1})\preceq_K K_H(\beta_r)
    \quad\text{and hence}\quad
    \mathcal{J}_H(\beta_{r+1};x)\le \mathcal{J}_H(\beta_r;x)
    \]
    for all $r\in\{0,\dots,R_{\mathrm{si}}-1\}$.
\end{enumerate}
\end{proposition}

\begin{proof}
By Definition~\ref{def:admissible-teacher}, $\mathcal{D}^{(0)}$ is an
admissible same-task teacher dataset. By construction, every record in each
$\mathcal{T}^{(r)}$ is a realized same-task exact label, so deterministic
merging preserves admissibility by induction over $r$. Hence part (1) follows:
each round is the same exact-$\mathcal{J}_H$ controller applied to a different
admissible warm start. Part (2) is immediate from the definition of
$\mathcal{T}^{(r)}$. For part (3), the sets
$\mathcal{S}_r^{\mathrm{meta}}$ are nested by construction, and $\beta_r$ is the
unique minimum of $K_H$ on $\mathcal{S}_r^{\mathrm{meta}}$. Enlarging the set
cannot worsen that minimum.
\end{proof}

\begin{theorem}[Conditional bounded self-improvement under predictor refinement]
\label{thm:aiqi-warmstart-self-improve}
Fix two rounds $r<s$ and a history $h_{<t}$ that occurs identically in both
rounds. Let $\psi^{(r)}$ and $\psi^{(s)}$ be the corresponding round-specific
unified phase predictors. If
\[
\sup_{a\in\mathcal{A}_B}
\left\|
\psi^{(s)}(\cdot\mid h_{<t}a)-p_t(\cdot\mid h_{<t},a)
\right\|_{\mathrm{TV}}
\le
\sup_{a\in\mathcal{A}_B}
\left\|
\psi^{(r)}(\cdot\mid h_{<t}a)-p_t(\cdot\mid h_{<t},a)
\right\|_{\mathrm{TV}},
\]
then at $h_{<t}$ the value-error bound and terminal-objective excess bound from
Theorem~\ref{thm:aiqi-warmstart-tv-robust} under round $s$ are no larger than
under round $r$. If the displayed inequality is strict, both bounds are strict.
\end{theorem}

\begin{proof}
Apply Theorem~\ref{thm:aiqi-warmstart-tv-robust} with
\[
\delta_t^{(r)}:=
\sup_{a\in\mathcal{A}_B}
\left\|
\psi^{(r)}(\cdot\mid h_{<t}a)-p_t(\cdot\mid h_{<t},a)
\right\|_{\mathrm{TV}},
\]
and the analogous definition of $\delta_t^{(s)}$. The bounds there are
$\overline G_m\delta_t$ for value error and $2\overline G_m\delta_t$ for
terminal-objective excess, so monotonicity and strictness follow immediately.
\end{proof}

\begin{corollary}[Advantage over bounded non-refresh reruns]
\label{cor:aiqi-warmstart-self-improve-vs-restart}
Consider a comparison procedure that reruns the same exact-$\mathcal{J}_H$
controller for the same bounded round deadlines but always reinitializes from
$\mathcal{D}^{(0)}$ instead of using $\mathcal{D}^{(r)}$. If for some later
round $s$ and some revisited history $h_{<t}$,
\[
\sup_{a\in\mathcal{A}_B}
\left\|
\psi^{(s)}(\cdot\mid h_{<t}a)-p_t(\cdot\mid h_{<t},a)
\right\|_{\mathrm{TV}}
<
\sup_{a\in\mathcal{A}_B}
\left\|
\psi^{(0)}(\cdot\mid h_{<t}a)-p_t(\cdot\mid h_{<t},a)
\right\|_{\mathrm{TV}},
\]
then the trace-refresh loop has a strictly smaller local finite-horizon
terminal-objective excess bound at $h_{<t}$ than the non-refresh rerun.
\end{corollary}

\begin{proof}
Apply Theorem~\ref{thm:aiqi-warmstart-self-improve} with round $r=0$. The
comparison procedure uses the same exact controller and the same objective;
only the predictor initialization differs.
\end{proof}

\begin{remark}[Self-improvement claim]
This is a bounded same-task self-improvement statement, not a universal
wall-clock dominance theorem. The exact claim is: every round is a real tuning
run on the original task, best-so-far deployable objective never worsens, and
later reruns can have strictly tighter same-state finite-horizon bounds
whenever accumulated real traces reduce predictor error on revisited histories.
\end{remark}

\subsection{Correspondence to Veness et al. and strict claim boundary}

The finite-horizon path above is exactly the structural move used for
$\rho$UCT consistency (Kocsis--Szepesv\'ari, 2006) as invoked by Veness et al.
(2010): interpret histories as Markov states, then apply finite-horizon MDP
consistency.
The explicit planner-convergence results in this document are
Theorem~\ref{thm:planner-convergence-mcaixi} and, for the broader set-valued
non-unique optimal-sequence case,
Proposition~\ref{prop:planner-convergence-optset}. The theorem's unique-path
assumption is a stronger sufficient hypothesis, not the main informal
takeaway; the assumptions are isolated in
Section~\ref{sec:planner-convergence-mcaixi}.

Any bounded-mixture or truncated-universal approximation used inside planner
model classes is a search-control/modeling choice only: deployability
predicate, evaluator semantics, incumbent key, and output definition remain
fixed by $(B,H)$ and $\mathcal{J}_H$.

The reset-augmented continuing path is the correct route when one additionally
wants stationary/ergodic finite-state Markov structure. Under injective
$\Psi_H$, this is the observed $n=1$ case of stationary/ergodic $n$-Markov
analysis.

Consistent with Veness et al. (2010), self-optimizing discussion for
MC-AIXI(FAC-CTW) carries an explicit rigor caveat: FAC-CTW with KT leaves is
not literally a countable Bayesian mixture, so the final self-optimizing claim
is supportive but not fully rigorous in that exact implementation model.

\begin{proposition}[Optimizer-invariant semantics]
Changing only the optimizer/controller (annealed hill climbing,
$\rho$UCT/MC-AIXI, discounted AIQI, or warm-start exact-$\mathcal{J}_H$
AIQI-style) can change traversal/search-control dynamics, but
cannot change candidate-validity rules, evaluator semantics, incumbent ordering
key, or final output definition once a trajectory is fixed.
\end{proposition}

\begin{proof}
These objects are defined independently of the optimizer transition/controller
law; changing the optimizer changes only which candidates are visited.
Scoring and output are then evaluated by the same fixed definitions on the
realized trajectory.
\end{proof}

\section{Conclusion}\label{sec:conclusion}

This specification defines a complete constrained tuning capability for
\texttt{infotheory}. It formalizes a single deployable two-part objective
$\mathcal{J}_H(z;x)=L_B(z)+C_x(z)$ with hard deployability constraints
(throughput and memory caps),
mutation grammar, annealed hill-climbing runtime,
MC-AIXI and AIQI-family runtime environment paths,
reward--observation interface, reproducibility requirements, and complexity
profile (including planner-internal simulation budget $N_{\mathrm{sim}}$). It is implementation-guiding for runtime-selectable controller
execution across annealed hill climbing, MC-AIXI(FAC-CTW),
the discounted AIQI controller and the warm-start exact-$\mathcal{J}_H$
controller, with the
explicit objective-aligned planner reward interface via
incumbent objective-difference reward. Finite-budget
guarantees are non-asymptotic (Section~\ref{sec:annealed-guarantees}); the
asymptotic theorems are assumption-conditional and out-of-budget scope:
Theorem~\ref{thm:idealized-asymptotic-hit} for annealed idealized dynamics and
Theorem~\ref{thm:planner-convergence-mcaixi} together with
Proposition~\ref{prop:planner-convergence-optset} for exact objective-aligned
undiscounted finite-horizon MC-AIXI planning, with the proposition capturing
the broader set-valued non-unique optimal-sequence case. The planner-convergence theorem
is specific to that MC-AIXI regime and is not a direct convergence theorem for
runtime discounted AIQI or runtime warm-start exact-$\mathcal{J}_H$ control
with finite simulations or exploratory/non-greedy control.
The MC-AIXI-warm-started model-free path is exact only in the local
finite-horizon sense above: with greedy deployment and an exact predictor, the
current action is optimal for the horizon-$m$ terminal-incumbent objective.
It is not the discounted AIQI theorem path; it is a new exact-$\mathcal{J}_H$
controller inspired by AIQI's delayed-label prediction construction.
The same exact mechanism also supports an optional bounded same-task
trace-refresh self-improvement loop: reruns remain on the original deployable
objective, reuse realized exact labels as teacher data, preserve monotone
best-so-far semantics, and can strictly tighten later local finite-horizon
bounds relative to bounded non-refresh reruns when predictor error decreases on
revisited histories.

Section~\ref{sec:restricted-mixture} adds the precise semantic layer: a
restricted Bayesian mixture over the bounded deployable semantic class, exact
deployable MDL/MAP equivalence with dominant-term codelength upper bound, and
an optional truncated-universality dominance baseline under the explicit
constant-overhead embedding assumption.

\end{document}
